\chapter{UVM Basics}
\label{ch:uvm}

\chapabstract{The Universal Verification Methodology is a class library
written in \SV{}. It is not an extension of the language, not a tool, and not
a new file format. In \cref{ch:tb} you built a transaction, a generator, a
driver, a monitor, a scoreboard, a coverage collector and an environment by
hand. This chapter puts a UVM base class in front of each of those six boxes,
shows what the library adds on top, and then builds a complete, runnable UVM
testbench for the CORDIC rotator. At the end you will be able to read an
industrial UVM testbench, write a small one, and judge honestly whether a
given project needs one at all.}

% A chapter-local convenience for the UVM macro names, which begin with a
% backtick.  \texttt{} would otherwise typeset the backtick as a left quote.

\section{What UVM is, and where it came from}
\label{sec:uvm:what}

UVM is a library of about 350 \SV{} classes, plus a few dozen macros, plus a
handful of package-scope functions. You compile it once, you \kw{import} it
like any other package, and you extend its classes like any other classes.
Everything it does could be done by hand, and in \cref{ch:tb} you did do it
by hand. Nothing in UVM required a change to IEEE 1800: every construct the
library is built from, classes and virtual methods, parameterised types,
queues, mailboxes and dynamic threads, is in \cref{ch:toolkit}. That is the
single most important sentence in this chapter, and if you keep it in mind the
rest of UVM stops being mysterious.

The point matters because of how UVM is usually taught. It arrives with its
own vocabulary (agents, sequencers, phases, objections, the factory), its own
naming convention, its own report format and its own command-line switches,
and the effect on a newcomer is that of meeting a new language. It is not.
\chFiveM{uvm\_info} is a macro that expands to an \kw{if} and a method call.
\code{uvm\_component} is a class with a \code{string} name and a parent
pointer. \code{run\_phase} is a virtual task that somebody else calls.

\begin{notebox}{Why the UVM listings are a different colour}
Throughout this chapter, listings that contain UVM identifiers use the
\code{uvmcode} environment. It typesets UVM class names, macros and phase
methods in rust, while the \SV{} keywords stay blue. Every rust identifier
you see is something defined in \file{uvm\_pkg.sv} or
\file{uvm\_macros.svh}. Nothing rust is part of IEEE 1800. Look at any
listing in this chapter and mentally delete the rust: what is left is the
language you already learned in \cref{ch:toolkit}.
\end{notebox}

\subsection{A short history}
\label{sec:uvm:history}

\index{UVM!history}%
Constrained-random, coverage-driven verification became mainstream around
2002, and each of the three big EDA vendors shipped a methodology library with
its simulator. Synopsys had VMM (Verification Methodology Manual), written
for VCS. Cadence had eRM, ported from the \emph{e} language into \SV{} as the
Universal Reuse Methodology, URM. Mentor had AVM (Advanced Verification
Methodology), which was the first of the three to be released under an open
licence and the first to be built around TLM communication rather than around
signals.

\index{OVM}\index{VMM}\index{AVM}%
In 2008 Mentor and Cadence merged AVM and URM into the Open Verification
Methodology, OVM. OVM was genuinely open: Apache-licensed source, downloadable,
and usable on any simulator. It was widely adopted, which left Synopsys with a
methodology nobody outside VCS used. In 2009 Accellera took OVM as the
starting point for a standard methodology, added the register layer from VMM
and a few ideas from both camps, and in February 2011 released UVM 1.0. The
version that most production code was written against is UVM~1.1d (2013); UVM
1.2 (2014) changed a few interfaces and was, for a while, unpopular precisely
because it broke code.

\index{IEEE 1800.2}%
In 2017 the IEEE standardised the class reference as IEEE 1800.2-2017,
revised in 2020. Note what was standardised: the \emph{API}, not the
implementation. Accellera continues to publish a reference implementation
(the current one identifies itself as UVM-1.2 / IEEE 1800.2-2020 depending on
the build), and every commercial simulator ships a copy, usually with vendor
patches for performance. You almost never download UVM yourself: you point
the compiler at the copy in the tool installation.

\subsection{What it costs}
\label{sec:uvm:cost}

Be honest with yourself about the price, because most of the UVM literature
is not.

\begin{itemize}
  \item \textbf{The learning curve is long.} A competent \SV{} programmer
        needs a week to write a first working UVM testbench and a month
        before the factory and the configuration database stop producing
        surprises. Compare that with \cref{ch:tb}, where the entire
        infrastructure was about 470 readable lines.
  \item \textbf{The code is verbose.} The minimal CORDIC testbench in
        \cref{sec:uvm:complete} is about 330 lines of UVM against about 470 lines
        of hand-written testbench, and the UVM version does less, because
        the hand-written one included the reference model and the reporting
        that UVM provides. Real testbenches are verbose in a different way:
        boiler-plate constructors, boiler-plate registrations,
        boiler-plate \code{build\_phase} bodies.
  \item \textbf{It is slow at run time.} The factory, the string-based
        configuration database, the field-automation macros and the objection
        mechanism all cost simulator time. On a small block the UVM overhead
        can dominate the \RTL{} simulation entirely. The field macros in
        particular are notorious: see \cref{sec:uvm:sequences}.
  \item \textbf{No open-source simulator runs it.} Be precise about why,
        because the usual summary is wrong. \tool{Verilator} compiles classes
        but not with the completeness UVM needs; it accepts a
        \kw{constraint} block and then ignores it, warning
        \code{CONSTRAINTIGN} and randomising as if the block were not there;
        and it rejects \kw{covergroup} outright as unsupported.
        \refsimulationbook explains what it does instead.
        \tool{Icarus Verilog} is often said not to support classes at all, and
        that is simply untrue: version 12 with \code{-g2012} compiles and runs
        a class with a constructor, member functions, inheritance and
        \code{super.new()}. What it lacks is everything built on top.
        \code{randomize()} does not exist (\code{randomize is not a method of
        class}) and \kw{constraint} blocks are rejected outright;
        \kw{covergroup} is a syntax error; concurrent
        assertions are rejected; and a queue of class handles is refused with
        \code{Sorry: Queue of type ... is not yet supported}, which alone rules
        out most of \file{uvm\_pkg.sv}. Even the class support that exists is
        fragile: two properties declared on one line, \code{int a, b;}, compile
        and then crash the runtime. Icarus is a useful Verilog-2005 simulator
        and a poor \SV{} one, and it will not run UVM.
\end{itemize}

In practice, then, UVM means a commercial licence: \tool{QuestaSim},
\tool{Xcelium} or \tool{VCS}, each of which ships a pre-compiled copy of the
library. For a student that means the university lab.

\begin{ruleofthumb}
UVM is a reuse technology. If nothing will be reused, not across blocks, not
across projects and not across people, you are paying the cost and collecting
none of the benefit.
\end{ruleofthumb}

\section{The same six boxes}
\label{sec:uvm:map}

Open \file{tb\_cordic\_pkg.sv} from \cref{ch:tb} next to this page. It
contains six classes: \code{cordic\_txn}, \code{cordic\_drv},
\code{cordic\_mon}, \code{cordic\_sb}, \code{cordic\_cov} and
\code{cordic\_env}. UVM has a base class for every one of them.
\cref{tab:uvm:map} is the whole translation.

\begin{table}[htbp]
  \centering
  \caption{Every class you wrote by hand in \cref{ch:tb}, its UVM base class,
           and what the base class actually adds.}
  \label{tab:uvm:map}
  \small
  \begin{tabularx}{\linewidth}{@{}l l X@{}}
    \toprule
    \textbf{\cref{ch:tb}} & \textbf{UVM base class} & \textbf{What UVM adds} \\
    \midrule
    \code{cordic\_txn} (\cref{ch:toolkit})
      & \code{uvm\_sequence\_item}
      & A name, a factory registration, automatic
        copy/compare/pack/print, and a handle back to the sequence that
        created it. \\
    \addlinespace
    \code{env.send\_random()}
      & \code{uvm\_sequence}
      & Stimulus becomes an object with its own class hierarchy: it can be
        randomised, parameterised, extended, nested inside another sequence,
        and selected at run time. \\
    \addlinespace
    \code{cordic\_drv}
      & \code{uvm\_driver \#(T)}
      & A standard blocking port to the sequencer
        (\code{seq\_item\_port}) instead of a bare mailbox, plus phasing. \\
    \addlinespace
    \code{cordic\_mon}
      & \code{uvm\_monitor}
      & Nothing, structurally. What you actually gain is the analysis port:
        one \code{write()} call feeds any number of subscribers, added
        without touching the monitor. \\
    \addlinespace
    \code{cordic\_sb}
      & \code{uvm\_scoreboard}
      & Again almost nothing by itself. The gain is
        \code{uvm\_tlm\_analysis\_fifo}, which decouples the monitor from the
        checking thread, and the report server, which counts your errors and
        decides the verdict. \\
    \addlinespace
    \code{cordic\_cov}
      & \code{uvm\_subscriber \#(T)}
      & A \code{write()} callback instead of a \code{forever} loop on a
        mailbox. One fewer process to start and stop. \\
    \addlinespace
    \code{cordic\_env}
      & \code{uvm\_env}
      & A named component tree, automatic phase execution, the configuration
        database and the factory. \\
    \addlinespace
    (the top-level module)
      & \code{uvm\_test}
      & The test becomes a class, chosen on the command line with
        \code{+UVM\_TESTNAME}. One compiled image, many tests. \\
    \addlinespace
    (none)
      & \code{uvm\_agent}
      & A packaging unit: sequencer + driver + monitor for one protocol, with
        an active/passive switch. This is the unit that gets reused. \\
    \addlinespace
    \code{mailbox \#(T)}
      & \code{uvm\_analysis\_port}, \code{uvm\_tlm\_analysis\_fifo}
      & One-to-many broadcast, a non-blocking producer, and connection made
        in \code{connect\_phase} rather than by passing handles into
        constructors. \\
    \addlinespace
    \code{wait\_drain()} + timeout
      & objections + drain time
      & End-of-test is negotiated by every component that has work
        outstanding, not decided by one counter in the test. \\
    \bottomrule
  \end{tabularx}
\end{table}

Read the right-hand column again. Half the rows say ``almost nothing''. That is
not a criticism of UVM; it is the point. UVM's value is not in any single
base class. It is in the fact that ten thousand engineers agree on the same
six boxes, the same names for them, the same connection mechanism and the
same start-up sequence, so that a verification component written by one team
drops into another team's testbench.

\subsection{The standard hierarchy}
\label{sec:uvm:hierarchy}

\cref{fig:uvm:hierarchy} is the picture every UVM testbench draws. Learn it
once and you will recognise the structure of any testbench you are handed.

\begin{figure}[htbp]
  \centering
  \begin{tikzpicture}[
      font=\sffamily\scriptsize,
      comp/.style={draw=codeframe,fill=svbg,line width=0.8pt,
                   rounded corners=2pt,minimum height=7mm,minimum width=20mm,
                   align=center},
      dutbox/.style={draw=codeframe,fill=vhdlbg,line width=0.8pt,
                     rounded corners=2pt,minimum height=8mm,minimum width=34mm,
                     align=center},
      frame/.style={draw=noteframe,dashed,line width=0.7pt,rounded corners=3pt},
      ana/.style={-{Stealth[length=2mm]},accentalt,line width=0.8pt},
      seqp/.style={{Stealth[length=2mm]}-{Stealth[length=2mm]},codekey,
                   line width=0.8pt},
      pin/.style={-{Stealth[length=2mm]},black!60,line width=0.7pt}
    ]

    % ---- agent internals -------------------------------------------------
    \node[comp] (sqr) at (0,0)      {sequencer};
    \node[comp] (drv) at (0,-1.5)   {driver};
    \node[comp] (mon) at (3.2,-1.5) {monitor};

    \node[frame,fit=(sqr)(drv)(mon),inner sep=4.5mm] (agent) {};
    \node[anchor=north west,font=\sffamily\scriptsize\bfseries,
          text=noteframe]
         at ([xshift=2pt,yshift=-1pt]agent.north west) {agent};

    % ---- analysis consumers ---------------------------------------------
    \node[comp,minimum width=24mm] (sb)  at (8.4,-0.6) {scoreboard};
    \node[comp,minimum width=24mm] (cov) at (8.4,-2.2) {coverage};

    \node[frame,fit=(agent)(sb)(cov),inner sep=4.5mm] (env) {};
    \node[anchor=north west,font=\sffamily\scriptsize\bfseries,
          text=noteframe]
         at ([xshift=2pt,yshift=-1pt]env.north west) {env};

    \node[comp,minimum width=26mm] (test) at (3.6,2.6) {test};

    % ---- DUT -------------------------------------------------------------
    \node[dutbox] (dut) at (1.6,-5.6) {DUT\\\texttt{cordic\_rot}};
    \node[draw=accent,fill=tipbg,line width=0.8pt,rounded corners=2pt,
          minimum width=52mm,minimum height=5mm,align=center]
         (ifc) at (1.6,-4.4) {\texttt{cordic\_if} (virtual interface)};

    % ---- edges -----------------------------------------------------------
    \draw[seqp] (sqr) -- node[right,font=\sffamily\tiny,align=left,xshift=1pt]
                          {seq\_item\_port} (drv);
    \draw[ana]  (mon.east) -- ++(1.7,0) |- (sb.west);
    \draw[ana]  (mon.east) -- ++(1.7,0) |- (cov.west);
    \node[font=\sffamily\tiny,text=accentalt,anchor=south west]
         at (4.35,-1.42) {analysis port};

    \draw[-{Stealth[length=2mm]},black!55,line width=0.7pt]
         (test.south) -- (env.north);

    \draw[pin] (drv.south) -- ++(0,-1.2) -| (-1.0,-4.15);
    \draw[{Stealth[length=2mm]}-] (mon.south) -- ++(0,-1.2) -| (3.2,-4.15);
    \draw[pin] (ifc.south) -- (dut.north);

  \end{tikzpicture}
  \caption{The standard UVM component hierarchy. Blue double arrows carry
    sequence items; rust arrows are analysis ports and are one-to-many.
    The driver and the monitor reach the pins through the same virtual
    interface, but only the driver ever writes to it.}
  \label{fig:uvm:hierarchy}
\end{figure}

Four things about that picture are worth stating explicitly, because they are
the ones students get wrong.

\begin{enumerate}
  \item \textbf{The test is at the top, and there is exactly one.} It is
        created by \code{run\_test()} from the class name given on the
        command line. Everything else is created, directly or indirectly, by
        the test's \code{build\_phase}.
  \item \textbf{The agent owns one protocol.} If the CORDIC had a
        configuration register bus as well as the data channel, it would get
        a second agent, with its own sequencer, driver and monitor.
  \item \textbf{The scoreboard and the coverage collector are outside the
        agent.} They are specific to the \DUT{}, not to the protocol. The
        agent is the reusable part; the scoreboard is not.
  \item \textbf{The arrows from the monitor are broadcast.} Adding a third
        subscriber, say a protocol checker, requires no change at all to the
        monitor. This is the concrete benefit over the two mailboxes you
        wired by hand in \cref{ch:tb}, where the monitor had to know that
        there were exactly two consumers.
\end{enumerate}

\section{\code{uvm\_object} versus \code{uvm\_component}}
\label{sec:uvm:obj-comp}

\index{UVM!uvm\_object}\index{UVM!uvm\_component}%
Before any of the mechanisms make sense, internalise this distinction. Every
class in UVM derives from \code{uvm\_void}, and immediately below it the tree
splits in two.

\begin{description}
  \item[\code{uvm\_object}] is \emph{data}. It has a name but no place. It is
        created, passed around, copied, compared and thrown away, thousands
        of times per simulation. Sequence items and sequences are objects.
        So are configuration objects and register-model fields.
  \item[\code{uvm\_component}] is \emph{structure}. It has a name
        \emph{and a parent}, so components form a tree. Components are
        created during \code{build\_phase}, they exist for the whole
        simulation, and they are never destroyed. Drivers, monitors,
        sequencers, agents, environments, scoreboards and tests are
        components.
\end{description}

The mental model that works is the \VHDL{} one. A \code{uvm\_component} is an
\emph{instance in an elaborated design hierarchy}: it corresponds to a
component instantiation inside an architecture, it has a path, and once
elaboration is over it is fixed. A \code{uvm\_object} is a
\emph{value of some type}: it corresponds to a variable of a record type
inside a process. \VHDL{} keeps those two worlds strictly apart, and so does
UVM.

\index{UVM!get\_full\_name}%
The consequence you will use every day is the hierarchical name. Because a
component knows its parent, \code{get\_full\_name()} returns a dotted path:

\begin{txtcode}
uvm_test_top.env.agt.drv
uvm_test_top.env.agt.mon
uvm_test_top.env.sb
\end{txtcode}

\code{uvm\_test\_top} is the name \code{run\_test()} gives to the test it
creates; everything below is what you named the handles in
\code{build\_phase}. Those strings are not cosmetic. They are the keys used
by the configuration database (\cref{sec:uvm:configdb}), by instance-specific
factory overrides (\cref{sec:uvm:factory}), and by verbosity control
(\cref{sec:uvm:reporting}). Every message printed by UVM carries the full
name of the component that produced it, which is how you find out which of
your four identical agents is the one misbehaving.

\begin{tipbox}{Name your handles after the instance, not the type}
Write \code{drv = cordic\_driver::type\_id::create("drv", this);}, not
\code{create("cordic\_driver", this)}. The string becomes the path element.
Short, consistent instance names such as \code{drv}, \code{mon}, \code{sqr},
\code{agt}, \code{env}, \code{sb} and \code{cov} make configuration paths
like \code{"env.agt.*"} readable. A common convention is to use the same
string as the handle name, so that the source and the printed topology look
the same.
\end{tipbox}

A \code{uvm\_object} has no parent and therefore no path. Its
\code{get\_full\_name()} returns just its own name, unless it is a sequence
running on a sequencer, in which case UVM splices in the sequencer's path.
This is why you cannot look up configuration ``for a transaction'': there is
nothing to look it up against.

\section{Phases}
\label{sec:uvm:phases}

\index{UVM!phases}%
In \cref{ch:tb} the environment's constructor built the sub-objects, then you
called \code{env.run()}, then you called \code{env.report()}. Construction,
connection, execution, reporting, in that order, by hand. UVM does the same
four things, calls them phases, and calls them for you.

A phase is a virtual method with a fixed name that UVM invokes on every
component in the tree. You override the ones you need and ignore the rest.

\begin{table}[htbp]
  \centering
  \caption{The common UVM phases, in execution order.}
  \label{tab:uvm:phases}
  \small
  \begin{tabularx}{\linewidth}{@{}l l l X@{}}
    \toprule
    \textbf{Phase} & \textbf{Kind} & \textbf{Order} & \textbf{What goes in it} \\
    \midrule
    \code{build\_phase}          & function & top-down
      & Create sub-components with the factory; fetch configuration. \\
    \code{connect\_phase}        & function & bottom-up
      & Connect TLM ports, exports and analysis ports. \\
    \code{end\_of\_elaboration\_phase} & function & bottom-up
      & Last chance to inspect the assembled tree; print the topology. \\
    \code{start\_of\_simulation\_phase} & function & bottom-up
      & Print a banner, display the resolved configuration. \\
    \code{run\_phase}            & \emph{task} & all in parallel
      & The only phase that consumes simulation time. \\
    \code{extract\_phase}        & function & bottom-up
      & Pull results out of components into a common place. \\
    \code{check\_phase}          & function & bottom-up
      & Check for leftovers: unmatched transactions, empty FIFOs. \\
    \code{report\_phase}         & function & bottom-up
      & Print your own summary. \\
    \code{final\_phase}          & function & top-down
      & Close files, flush databases. \\
    \bottomrule
  \end{tabularx}
\end{table}

\begin{figure}[htbp]
  \centering
  \begin{tikzpicture}[
      font=\sffamily\scriptsize,
      ph/.style={draw=codeframe,fill=svbg,line width=0.8pt,rounded corners=2pt,
                 minimum width=46mm,minimum height=6.5mm,align=center},
      run/.style={draw=accentalt,fill=warnbg,line width=1pt,rounded corners=2pt,
                  minimum width=46mm,minimum height=6.5mm,align=center},
      fl/.style={-{Stealth[length=2mm]},black!60,line width=0.8pt},
      dir/.style={-{Stealth[length=1.6mm]},codekey,line width=0.7pt}
    ]
    \node[ph] (b)  at (0, 0.0)  {\texttt{build\_phase}};
    \node[ph] (c)  at (0,-1.0)  {\texttt{connect\_phase}};
    \node[ph] (e)  at (0,-2.0)  {\texttt{end\_of\_elaboration\_phase}};
    \node[ph] (s)  at (0,-3.0)  {\texttt{start\_of\_simulation\_phase}};
    \node[run](r)  at (0,-4.3)  {\texttt{run\_phase}\ \ (task, time passes)};
    \node[ph] (x)  at (0,-5.6)  {\texttt{extract\_phase}};
    \node[ph] (k)  at (0,-6.6)  {\texttt{check\_phase}};
    \node[ph] (rp) at (0,-7.6)  {\texttt{report\_phase}};
    \node[ph] (f)  at (0,-8.6)  {\texttt{final\_phase}};

    \foreach \a/\b in {b/c,c/e,e/s,s/r,r/x,x/k,k/rp,rp/f}
      \draw[fl] (\a) -- (\b);

    % direction annotations
    \draw[dir] (3.0,0.45) -- (3.0,-0.45);
    \node[anchor=west,font=\sffamily\tiny,text=codekey] at (3.15,0.0)
         {top-down: parent creates children};
    \draw[dir] (3.0,-1.45) -- (3.0,-0.55);
    \node[anchor=west,font=\sffamily\tiny,text=codekey] at (3.15,-1.0)
         {bottom-up: children exist first};
    \draw[dir] (3.0,-2.45) -- (3.0,-1.55);
    \draw[dir] (3.0,-3.45) -- (3.0,-2.55);
    \node[anchor=west,font=\sffamily\tiny,text=codekey] at (3.15,-3.0)
         {bottom-up};
    \node[anchor=west,font=\sffamily\tiny,text=accentalt] at (3.15,-4.3)
         {every component's \texttt{run\_phase} forked in parallel;};
    \node[anchor=west,font=\sffamily\tiny,text=accentalt] at (3.15,-4.75)
         {ends when the last objection is dropped};
    \draw[dir] (3.0,-6.05) -- (3.0,-5.15);
    \node[anchor=west,font=\sffamily\tiny,text=codekey] at (3.15,-6.6)
         {bottom-up};
    \draw[dir] (3.0,-7.05) -- (3.0,-6.15);
    \draw[dir] (3.0,-8.05) -- (3.0,-7.15);
    \draw[dir] (3.0,-8.25) -- (3.0,-9.05);
    \node[anchor=west,font=\sffamily\tiny,text=codekey] at (3.15,-8.6)
         {top-down};

    % run-time sub-phases
    \node[draw=codeframe,fill=codebg,line width=0.6pt,rounded corners=2pt,
          minimum width=46mm,minimum height=5mm,font=\sffamily\tiny,align=center]
         (sub) at (0,-10.2)
         {optional run-time sub-phases, in parallel with \texttt{run\_phase}:\\
          \texttt{reset} $\rightarrow$ \texttt{configure} $\rightarrow$
          \texttt{main} $\rightarrow$ \texttt{shutdown}};
    \draw[fl,dashed] (r.west) -- ++(-1.1,0) |- (sub.west);
  \end{tikzpicture}
  \caption{UVM phase order. Only \code{run\_phase} consumes simulation time;
    everything above and below it happens at a single time step.}
  \label{fig:uvm:phases}
\end{figure}

\subsection{Why build is top-down and connect is bottom-up}
\label{sec:uvm:phase-order}

This is not an arbitrary choice, and understanding it removes most
build-order confusion.

\code{build\_phase} must be top-down because a parent \emph{creates} its
children. The test's \code{build\_phase} creates the environment; only then
does the environment exist, so only then can UVM call the environment's
\code{build\_phase}, which creates the agent, and so on. The tree grows
downwards as the phase executes. A useful side effect: because the parent
runs first, the parent can write configuration into the database that the
child will read a moment later in its own \code{build\_phase}. That is
exactly how the active/passive setting and the virtual interface reach the
driver.

\code{connect\_phase} must be bottom-up because a connection needs both ends
to exist. When the agent connects \code{drv.seq\_item\_port} to
\code{sqr.seq\_item\_export}, both the driver and the sequencer must already
have been constructed \emph{and} must already have constructed their own
ports. Running the deepest components first guarantees it.

The remaining function phases are bottom-up for the same reason in reverse:
by the time the environment reports, its children have already reported, so
the environment can summarise them. \code{final\_phase} is top-down only for
symmetry with \code{build\_phase}.

\subsection{\code{super.build\_phase(phase)}}
\label{sec:uvm:super-build}

Every phase method you override should begin by calling the base-class
version:

\begin{uvmcode}[caption={The shape of every component},label={lst:uvm:shape}]
class cordic_agent extends uvm_agent;
  `uvm_component_utils(cordic_agent)

  function new(string name, uvm_component parent);
    super.new(name, parent);
  endfunction

  function void build_phase(uvm_phase phase);
    super.build_phase(phase);     // never omit this line
    // ... create children here ...
  endfunction
endclass
\end{uvmcode}

\begin{gotcha}{Omitting \code{super.build\_phase(phase)}}
\index{UVM!build\_phase}%
This is the single most common UVM bug, and its symptom is baffling: things
silently do not get configured.

\code{uvm\_component::build\_phase} is not empty. It calls
\code{apply\_config\_settings()}, which walks the configuration database and
automatically assigns any registered field of this component whose name
matches a setting. \code{uvm\_agent::build\_phase} additionally reads the
\code{is\_active} setting. If you skip \code{super.build\_phase(phase)},
none of that happens: your agent stays \code{UVM\_ACTIVE} regardless of what
the test asked for, and any field you expected the automation to fill stays
at its default.

The failure produces no error and no warning. It produces an agent that
drives pins when it was supposed to be passive.

In \VHDL{} terms: you overrode a procedure in a derived type and forgot that
the parent's version was doing real work. \VHDL{} has no inheritance, so
nothing in your previous experience warns you about this.
\end{gotcha}

\subsection{The run-time sub-phases}
\label{sec:uvm:subphases}

\index{UVM!run-time sub-phases}%
UVM also defines a second schedule of run-time sub-phases that executes in
parallel with \code{run\_phase}. There are four of them, namely
\code{reset}, \code{configure}, \code{main} and \code{shutdown}, and
each has a \code{pre\_} and a \code{post\_} variant, twelve in all. Each is
a task, each is objected to independently, and the intent is that a testbench can, for example, apply
reset in \code{reset\_phase} and stimulus in \code{main\_phase} without the
two getting tangled.

\begin{notebox}{Nobody uses the sub-phases}
That is an exaggeration, but only a small one. Many teams -- probably most --
put everything in \code{run\_phase} and manage ordering with ordinary
\kw{fork}/\kw{join} and events. The sub-phases add a second synchronisation
mechanism on top of objections, they interact badly with sequences that span
phases, and \emph{jumping} between them (\code{phase.jump()}) is a source of
hard-to-debug behaviour. If you are learning UVM, use \code{run\_phase} only.
If you join a team that uses the sub-phases, follow the local convention.
\end{notebox}

\section{Objections: how a test knows it is finished}
\label{sec:uvm:objections}

\index{UVM!objection}%
In \cref{ch:tb} the test finished because you wrote \code{wait\_drain()}: a
loop that polled the scoreboard's transaction count against an expected
number, with a cycle-count timeout that called \code{\$fatal} if the count
never arrived. That works, but it puts the knowledge of ``am I done'' in one
place, the test, and that place has to know how many transactions everyone
else is still holding.

UVM inverts it. Any component that has outstanding work \emph{objects} to the
phase ending. The phase ends when the objection count reaches zero and stays
zero for the drain time. Nobody needs to know the global count.

\begin{uvmcode}[caption={An objection around the stimulus},label={lst:uvm:objection}]
task run_phase(uvm_phase phase);
  cordic_random_seq seq;

  phase.raise_objection(this, "stimulus running");
  phase.phase_done.set_drain_time(this, 500ns);

  seq = cordic_random_seq::type_id::create("seq");
  seq.start(env.agt.sqr);

  phase.drop_objection(this, "stimulus done");
endtask
\end{uvmcode}

\index{UVM!drain time}%
The drain time is the direct equivalent of the margin you built into
\code{wait\_drain}'s timeout, but it means the opposite thing. It is not a
limit: it is a grace period. When the last objection drops, UVM waits 500\,ns
of simulation time before killing the run phase, which gives the pipeline in
the \DUT{} time to push out the last results, the monitor time to see them,
and the scoreboard time to check them. Without it, the test would end at the
instant the last request was accepted and the last few responses would never
be checked. Five hundred nanoseconds at a 10\,ns clock is 50 cycles, which
comfortably covers the CORDIC's $N+2$ cycle latency.

The second argument to \code{raise\_objection} and \code{drop\_objection} is
a description string. It costs nothing and it is invaluable:
\code{+UVM\_OBJECTION\_TRACE} on the command line prints every raise and drop
with its description, which is how you find out who is refusing to let your
test finish.

\subsection{Where to put them}

\begin{ruleofthumb}
Raise and drop objections in the test, around the sequence start. Everywhere
else, prefer not to.
\end{ruleofthumb}

You \emph{can} raise an objection in a driver, a monitor or a sequence, and
sometimes you must, since a monitor still holding a half-observed transaction
has a legitimate claim that the test is not over. But every
raise/drop pair is a place where the test can hang, and objections raised
deep in the tree are hard to trace. Start with one pair, in the test.

\begin{gotcha}{A test that ends at time zero}
\index{UVM!objection}%
You run your brand-new testbench and the transcript says:

\begin{txtcode}
UVM_INFO ... [UVM/REPORT/SERVER]
--- UVM Report Summary ---
** Report counts by severity
UVM_INFO :    9
UVM_WARNING : 0
UVM_ERROR :   0
UVM_FATAL :   0
$finish called from file "uvm_root.svh"
\end{txtcode}

at simulation time 0, with no transactions and no waveform. Nobody raised an
objection, so the run phase had nothing to wait for and ended immediately.

The \emph{opposite} bug -- a test that never ends -- has the same cause
inverted: an objection raised and never dropped, usually because the
\code{drop\_objection} sits after a line that blocks forever. Run with
\code{+UVM\_OBJECTION\_TRACE} and read the last raise.

Note that a component whose \code{run\_phase} is a plain \kw{forever} loop
(the driver, the monitor) does \emph{not} keep the phase alive. UVM kills
those threads when the phase ends. Only objections keep a phase alive.
\end{gotcha}

\begin{tipbox}{Always set a global timeout}
An objection that is never dropped hangs the simulation. Bound it:

\begin{uvmcode}
function void build_phase(uvm_phase phase);
  super.build_phase(phase);
  uvm_top.set_timeout(1ms, 1);   // hard stop, overridable
  env = cordic_env::type_id::create("env", this);
endfunction
\end{uvmcode}

or pass \code{+UVM\_TIMEOUT=1ms,YES} on the command line. A regression that
hangs costs a licence for hours; a regression that fails takes a minute.
\end{tipbox}

\section{The factory}
\label{sec:uvm:factory}

\index{UVM!factory}%
The factory is the mechanism that most repays understanding, because it is
the reason UVM testbenches are reusable at all.

\subsection{The problem it solves}

Suppose you have a working environment, and you now want a test in which the
driver occasionally corrupts a transaction, so that you can check that the
scoreboard notices. In \cref{ch:tb} you would have to edit
\code{cordic\_env::new()}, because that is where \code{drv = new(...)} is
written. Editing a working, shared environment in order to add one test is
exactly what you must not do.

The factory removes the type name from the point of instantiation. Instead of
naming the class you want, you ask a registry for ``whatever class is
currently registered under this type'', and a test can change the answer
without touching the environment.

\begin{uvmcode}[caption={Registration and creation},label={lst:uvm:factory-basics}]
class cordic_driver extends uvm_driver #(cordic_item);
  `uvm_component_utils(cordic_driver)      // registers the type
  // ...
endclass

class cordic_item extends uvm_sequence_item;
  `uvm_object_utils(cordic_item)           // registers the type
  // ...
endclass

// inside cordic_agent::build_phase
drv = cordic_driver::type_id::create("drv", this);   // ask the factory
// NOT:  drv = new("drv", this);
\end{uvmcode}

\chFiveM{uvm\_component\_utils} and \chFiveM{uvm\_object\_utils} expand to a
typedef of a parameterised registry class (\code{type\_id}), a static
\code{get\_type()} function, a \code{create()} function and a
\code{get\_type\_name()} override, and they run a static registration that
inserts the type into the global factory. \code{type\_id::create()} then asks
the factory for an instance.

\subsection{Overriding a type}

\index{UVM!factory override}%
With registration in place, a new test can substitute a derived class for the
original everywhere it is created:

\begin{uvmcode}[caption={An error-injecting driver, substituted by a test},label={lst:uvm:override}]
// -------- the new driver: it exists only in the test's own file --------
class cordic_bad_driver extends cordic_driver;
  `uvm_component_utils(cordic_bad_driver)

  int unsigned corrupt_one_in = 20;

  function new(string name, uvm_component parent);
    super.new(name, parent);
  endfunction

  // Same signature as the parent's.  The parent declares drive() as
  // `virtual', so the parent's run_phase dispatches to this one.
  protected virtual task drive(cordic_item it);
    if ($urandom_range(corrupt_one_in - 1) == 0) begin
      `uvm_info(get_type_name(),
                $sformatf("injecting error into %s", it.convert2string()),
                UVM_LOW)
      it.theta = it.theta ^ 16'h0040;      // flip one bit
    end
    super.drive(it);
  endtask
endclass

// -------- the test: no change at all to the env or the agent -----------
class cordic_inject_test extends cordic_base_test;
  `uvm_component_utils(cordic_inject_test)

  function new(string name, uvm_component parent);
    super.new(name, parent);
  endfunction

  function void build_phase(uvm_phase phase);
    // The override MUST be set before the component is created, so it goes
    // before super.build_phase(), which is what creates the env.
    set_type_override_by_type(cordic_driver::get_type(),
                              cordic_bad_driver::get_type());
    super.build_phase(phase);
  endfunction

  task run_phase(uvm_phase phase);
    cordic_random_seq seq;
    phase.raise_objection(this);
    seq = cordic_random_seq::type_id::create("seq");
    void'(seq.randomize() with { n_items == 200; });
    seq.start(env.agt.sqr);
    phase.drop_objection(this);
  endtask
endclass
\end{uvmcode}

Run this test and every \code{cordic\_driver} in the tree becomes a
\code{cordic\_bad\_driver}. The environment source file is untouched. The
agent source file is untouched. Nothing was recompiled that did not have to
be.

If there are several agents and you want only one of them corrupted, use an
instance override, which takes a path string:

\begin{uvmcode}
set_inst_override_by_type("env.agt_a.drv",
                          cordic_driver::get_type(),
                          cordic_bad_driver::get_type());
\end{uvmcode}

Instance overrides take priority over type overrides, and the path is matched
against \code{get\_full\_name()} relative to the component that sets it. There
are also \code{\_by\_name} variants of both calls, which take a type name as
a string instead of a type handle; prefer the \code{\_by\_type} forms, because a typo in a type
name is a silent no-op while a typo in a type handle is a compile error.

\subsection{This is \VHDL{}'s \kw{configuration}, done with classes}

\index{configuration (VHDL)}%
If you have used \VHDL{}'s configuration declarations, the factory will feel
familiar, because they solve the same problem: choose which implementation
fills a hole in the hierarchy, without editing the hierarchy.

\begin{rosetta}
\begin{rleft}
-- one architecture is written
-- against a component declaration
architecture struct of top is
  component drv_c is
    port (...);
  end component;
begin
  u_drv : drv_c port map (...);
end architecture;

-- a separate configuration unit
-- decides what fills the hole
configuration cfg_bad of top is
  for struct
    for u_drv : drv_c
      use entity work.drv_bad(rtl);
    end for;
  end for;
end configuration;
\end{rleft}
\begin{rright}
// the agent is written against
// the base type
drv = cordic_driver::type_id
        ::create("drv", this);

// a separate test decides what
// fills the hole
class cordic_inject_test
  extends cordic_base_test;
  function void build_phase(
      uvm_phase phase);
    set_type_override_by_type(
      cordic_driver::get_type(),
      cordic_bad_driver::get_type());
    super.build_phase(phase);
  endfunction
endclass
\end{rright}
\end{rosetta}

The differences are instructive. The \VHDL{} configuration is elaborated
once, statically, and selects an entity/architecture pair. The UVM override
is executed at run time by ordinary code, so it can be conditional, can be
driven by a plusarg, and can be undone. On the other hand the \VHDL{}
configuration is checked by the analyser, while a UVM override of an
unrelated type fails only when the factory tries the \code{\$cast} and
reports an error at run time.

\begin{gotcha}{\kw{new} instead of \code{create}}
\index{UVM!type\_id::create}%
\begin{uvmcode}
drv = new("drv", this);                            // legal, and wrong
drv = cordic_driver::type_id::create("drv", this); // right
\end{uvmcode}

The first line compiles, runs, and produces a perfectly functional driver.
It simply bypasses the factory, so every override of that type silently does
nothing. There is no warning. Your error-injection test will run happily and
inject nothing.

The rule has no exceptions worth learning as a beginner: inside a testbench,
components and sequence items are always built with
\code{type\_id::create()}. The only \kw{new} calls in a UVM testbench are in
constructors (\code{super.new(name, parent)}) and for objects that are
genuinely not factory-registered, such as TLM ports and analysis ports.
\end{gotcha}

\begin{tipbox}{Check what the factory actually did}
Call \code{uvm\_factory::get().print()} in
\code{end\_of\_elaboration\_phase}, or pass \code{+UVM\_CONFIG\_DB\_TRACE}
and \code{-uvmtrace} depending on your simulator. The factory print lists
every registered type and every override in force. If your override is not in
that list, it was never applied.
\end{tipbox}

\section{The configuration database}
\label{sec:uvm:configdb}

\index{UVM!uvm\_config\_db}%
The factory decides \emph{what type} is built. The configuration database
decides \emph{what values} it is built with. It is a global, string-keyed,
type-parameterised store:

\begin{uvmcode}
uvm_config_db #(T)::set(uvm_component cntxt, string inst_name,
                        string field_name, T value);

bit ok = uvm_config_db #(T)::get(uvm_component cntxt, string inst_name,
                                 string field_name, ref T value);
\end{uvmcode}

The lookup key is the concatenation of \code{cntxt.get\_full\_name()} and
\code{inst\_name}, so the two arguments together name a scope. Passing
\kw{null} as the context means ``the root'', and then \code{inst\_name} is an
absolute path from the top. The path may contain \code{*} and \code{?}
wildcards, and when several settings match, the one whose scope is
\emph{higher} in the tree wins, with later settings beating earlier ones at
the same level. That precedence rule is what lets a test override what an
environment set.

\subsection{Getting the virtual interface into the driver}
\label{sec:uvm:vif}

\index{virtual interface}%
This is the one use of the configuration database that every testbench needs,
and it is worth writing out in full because the two halves live in different
files.

The \DUT{} and the physical interface are in a module. Classes cannot be
instantiated inside a module hierarchy, and a class cannot reach into a
module by name in any portable way, so the handle has to be passed. The
top-level module puts it into the database before \code{run\_test()} returns
control to the phasing:

\begin{uvmcode}[caption={Top-level module: publish the interface handle},label={lst:uvm:vif-set}]
module tb_cordic_uvm;
  import uvm_pkg::*;
  `include "uvm_macros.svh"
  import cordic_pkg::*;
  import cordic_uvm_pkg::*;

  logic clk = 1'b0;
  always #5ns clk = ~clk;

  cordic_if vif (.clk(clk));

  cordic_rot #(.ITER(N)) dut ( /* ports: see the full listing below */ );

  initial begin
    uvm_config_db #(virtual cordic_if)::set(null, "*", "vif", vif);
    run_test();
  end
endmodule
\end{uvmcode}

and every component that needs pins fetches it in its \code{build\_phase}:

\begin{uvmcode}[caption={Driver and monitor: fetch the interface handle},label={lst:uvm:vif-get}]
function void build_phase(uvm_phase phase);
  super.build_phase(phase);
  if (!uvm_config_db #(virtual cordic_if)::get(this, "", "vif", vif))
    `uvm_fatal("NOVIF",
               {"no virtual interface set for ", get_full_name(), ".vif"})
endfunction
\end{uvmcode}

Read the arguments carefully.

\begin{itemize}
  \item In the \code{set}, the context is \kw{null} and the instance name is
        \code{"*"}: available to every component in the tree. In a testbench
        with two interfaces you would write \code{"*.agt\_a.*"} and
        \code{"*.agt\_b.*"} instead of \code{"*"}.
  \item In the \code{get}, the context is \kw{this} and the instance name is
        the empty string: look up the key exactly for this component. Almost
        every \code{get} you will ever write has this form.
  \item The field name \code{"vif"} is a string and must match on both sides.
        It is conventional, not required, that it equals the name of the
        class member it is assigned to.
  \item The type parameter must match exactly. \code{virtual cordic\_if} and
        \code{virtual cordic\_if.drv} (a modport-qualified virtual interface)
        are \emph{different} types and will not find each other.
\end{itemize}

\begin{gotcha}{The null virtual interface}
Forget the \code{set}, misspell the field name, or use a wildcard that does
not match, and the \code{get} returns 0. If you did not test the return
value, \code{vif} stays \kw{null}, and the first clocking-block access in
\code{run\_phase} produces something like

\begin{txtcode}
** Fatal: (vsim-3363) Null virtual interface reference.
   Time: 0 ns  Iteration: 0  Process:
   /tb_cordic_uvm/#INITIAL#33/uvm_pkg::run_test
\end{txtcode}

or, on other simulators,

\begin{txtcode}
Error-[NOA] Null object access
  The object is being used before it was constructed/allocated.
\end{txtcode}

Neither message mentions the configuration database, the field name, or the
component. This is why the \chFiveM{uvm\_fatal} guard in
\cref{lst:uvm:vif-get} is not optional style: it converts an unreadable
simulator error at an arbitrary time into a message at time zero that names
the component and the missing key.

The most common cause among students: the \code{set} was written
\emph{after} \code{run\_test()} in the same \kw{initial} block.
\code{run\_test()} does not return until the simulation is over.
\end{gotcha}

\begin{notebox}{\code{uvm\_config\_db} versus \code{uvm\_resource\_db}}
\code{uvm\_config\_db} is a thin, hierarchy-aware layer over
\code{uvm\_resource\_db}. The resource database has no notion of scope
precedence; the configuration database adds the ``higher in the tree wins''
rule. Use \code{uvm\_config\_db} unless you have a specific reason not to.
Both are string-keyed and both fail silently on a typo, so use
\code{+UVM\_CONFIG\_DB\_TRACE} when a value does not arrive.
\end{notebox}

\begin{tipbox}{Pass a configuration object, not fifteen fields}
Once a component needs more than two or three settings, stop putting
individual fields in the database. Define a class:

\begin{uvmcode}
class cordic_cfg extends uvm_object;
  `uvm_object_utils(cordic_cfg)
  int unsigned          stall_pct  = 25;
  int unsigned          tol_lsb    = 12;
  uvm_active_passive_enum is_active = UVM_ACTIVE;
  function new(string name = "cordic_cfg"); super.new(name); endfunction
endclass
\end{uvmcode}

and put one handle in the database:
\code{uvm\_config\_db \#(cordic\_cfg)::set(this, "env.agt*", "cfg", cfg);}.
One key, one lookup, and the settings are type-checked by the compiler
instead of by a string comparison.
\end{tipbox}

\section{Sequences and sequencers}
\label{sec:uvm:sequences}

\index{UVM!sequence}\index{UVM!sequencer}%
In \cref{ch:tb} the stimulus was a task on the environment:
\code{send\_random(n)} looped $n$ times, created a transaction, randomised it
and pushed it into a mailbox. That is adequate and it is easy to read, but it
has a structural weakness: the stimulus is a \emph{task}, and a task cannot
be extended, cannot be randomised as a whole, and cannot be selected by name
at run time. UVM makes stimulus an \emph{object} instead, and that single
change is what makes test-writing scale.

\subsection{The sequence item}

A sequence item is a \code{uvm\_object} that carries one transaction. It is
the UVM version of \code{cordic\_txn}.

\begin{uvmcode}[caption={A sequence item with field automation},label={lst:uvm:item-macros}]
class cordic_item extends uvm_sequence_item;

  rand data_t  x0;
  rand data_t  y0;
  rand angle_t theta;
  data_t       x, y;        // filled in by the monitor

  `uvm_object_utils_begin(cordic_item)
    `uvm_field_int(x0,    UVM_DEFAULT | UVM_DEC)
    `uvm_field_int(y0,    UVM_DEFAULT | UVM_DEC)
    `uvm_field_int(theta, UVM_DEFAULT | UVM_DEC)
    `uvm_field_int(x,     UVM_DEFAULT | UVM_DEC)
    `uvm_field_int(y,     UVM_DEFAULT | UVM_DEC)
  `uvm_object_utils_end

  function new(string name = "cordic_item");
    super.new(name);
  endfunction

endclass
\end{uvmcode}

\index{UVM!field automation macros}%
The block between \chFiveM{uvm\_object\_utils\_begin} and
\chFiveM{uvm\_object\_utils\_end} is the \emph{field automation}. Each
\chFiveM{uvm\_field\_int} adds the named field to a table that the base class
walks whenever somebody calls \code{copy()}, \code{compare()}, \code{print()},
\code{pack()}, \code{unpack()}, \code{record()} or \code{sprint()}. You get
seven methods for five lines. The flags select which of them the field takes
part in: \code{UVM\_DEFAULT} means all of them, \code{UVM\_DEC} is a
formatting hint, and you can subtract behaviours with, for example,
\code{UVM\_DEFAULT | UVM\_NOCOMPARE} for a field the scoreboard should
ignore.

\begin{notebox}{Field macros are convenient and slow}
\index{UVM!field automation macros}%
The automation is implemented by a generic recursive walk over a table of
field descriptors, driven by string comparisons and \code{\$cast}. It is one
to two orders of magnitude slower than the hand-written equivalent, and it is
executed on every single transaction. On a testbench that moves millions of
items, the field macros can become a visible fraction of total run time.

Many experienced teams therefore do not use them at all. They register the
type with the plain \chFiveM{uvm\_object\_utils} and write the three methods
they actually need:

\begin{uvmcode}
`uvm_object_utils(cordic_item)

function void do_copy(uvm_object rhs);
  cordic_item that;
  super.do_copy(rhs);
  if (!$cast(that, rhs))
    `uvm_fatal("CAST", "do_copy: type mismatch")
  x0 = that.x0;  y0 = that.y0;  theta = that.theta;
  x  = that.x;   y  = that.y;
endfunction

function bit do_compare(uvm_object rhs, uvm_comparer comparer);
  cordic_item that;
  if (!$cast(that, rhs)) return 0;
  return super.do_compare(rhs, comparer)
         && (x == that.x) && (y == that.y);
endfunction

function string convert2string();
  return $sformatf("x0=%0d y0=%0d theta=%0d -> x=%0d y=%0d",
                   x0, y0, theta, x, y);
endfunction
\end{uvmcode}

That is more lines, but it is explicit, fast, and it is what you will see in
most industrial code. Use the macros while you are learning; be ready to
delete them when a profiler tells you to.
\end{notebox}

\begin{gotcha}{\chFiveM{uvm\_field\_int} on a \kw{real}}
\index{UVM!field automation macros}%
The field macros come in typed flavours and you must pick the right one:
\chFiveM{uvm\_field\_int} for integral types, \chFiveM{uvm\_field\_real} for
\kw{real} and \kw{shortreal}, \chFiveM{uvm\_field\_enum},
\chFiveM{uvm\_field\_string}, \chFiveM{uvm\_field\_object},
\chFiveM{uvm\_field\_array\_int}, \chFiveM{uvm\_field\_queue\_object}, and so
on, roughly thirty of them.

Writing \chFiveM{uvm\_field\_int(gain, UVM\_DEFAULT)} where \code{gain} is a
\kw{real} does not always fail to compile. The macro body assigns through a
\code{uvm\_bitstream\_t} (a 4096-bit vector), so a \kw{real} may be implicitly
converted to an integer, and the fractional part is silently lost. Your
comparison then passes on values that differ, or your print shows
\code{1} where the value is 1.6468.

There is a deeper reason not to put a \kw{real} in a transaction at all: a
transaction should hold what crossed the pins, and pins carry bits. Keep the
fixed-point value; convert to \kw{real} inside the scoreboard, where the
reference model lives.
\end{gotcha}

\subsection{The sequence}

A sequence is a \code{uvm\_object} whose \code{body()} task generates items.
It is the direct replacement for \code{send\_random()}.

\begin{uvmcode}[caption={A random sequence},label={lst:uvm:seq-basic}]
class cordic_random_seq extends uvm_sequence #(cordic_item);
  `uvm_object_utils(cordic_random_seq)

  rand int unsigned n_items = 200;
  constraint c_n { n_items inside {[1:5000]}; }

  function new(string name = "cordic_random_seq");
    super.new(name);
  endfunction

  task body();
    cordic_item it;
    repeat (n_items) begin
      it = cordic_item::type_id::create("it");
      start_item(it);
      if (!it.randomize())
        `uvm_fatal(get_type_name(), "item randomization failed")
      finish_item(it);
    end
  endtask

endclass
\end{uvmcode}

Three details deserve attention.

\textbf{The parameter.} \code{uvm\_sequence \#(cordic\_item)} is really
\code{uvm\_sequence \#(REQ, RSP)} with \code{RSP} defaulting to \code{REQ}.
Unless you use the response path (\cref{sec:uvm:pipelined}), one parameter is
enough.

\textbf{The randomisation is late.} The item is randomised \emph{between}
\code{start\_item} and \code{finish\_item}, not before \code{start\_item}.
That is deliberate. \code{start\_item} blocks until the sequencer grants this
sequence the right to send, and only then does the item become the ``current
item''; randomising after the grant means the constraints see the state of the
world at the moment the item will actually be driven, and it allows a
sequencer-level constraint layer to influence the result. Randomising before
\code{start\_item} works too, and is common, but it is the less general
pattern.

\textbf{The sequence itself is randomisable.} \code{n\_items} is \kw{rand},
so the test can write
\code{seq.randomize() with \{ n\_items inside \{[100:200]\}; \}}. This is the
capability a task simply does not have.

The \chFiveM{uvm\_do} family compresses the create/start/randomise/finish
sequence into one line:

\begin{uvmcode}
task body();
  cordic_item it;
  repeat (n_items)
    `uvm_do_with(it, { theta inside {[-PI_EXT : PI_EXT]}; })
endtask
\end{uvmcode}

\chFiveM{uvm\_do}, \chFiveM{uvm\_do\_with}, \chFiveM{uvm\_create},
\chFiveM{uvm\_send} and their \code{\_on} variants (which take an explicit
sequencer) are still widely used. They are also widely disliked: they hide
the handshake you are trying to learn, they make the constraint block hard to
read, and their error messages point inside the macro. Write the long form
until the handshake is second nature.

\subsection{Starting a sequence}

A sequence does not run by itself. Somebody calls \code{start()} on it and
tells it which sequencer to run on:

\begin{uvmcode}
cordic_random_seq seq;
seq = cordic_random_seq::type_id::create("seq");
void'(seq.randomize() with { n_items == 500; });
seq.start(env.agt.sqr);      // blocks until body() returns
\end{uvmcode}

\code{start()} is a task and it blocks, which is exactly what you want inside
an objection. The full signature is
\code{start(sequencer, parent\_seq, priority, call\_pre\_post)}; the extra
arguments matter only for nested sequences.

Inside \code{body()}, the handle \code{m\_sequencer} refers to the sequencer
the sequence is running on, typed as \code{uvm\_sequencer\_base}. If your
sequencer is a derived class with extra members, say a pointer to a
configuration object or a handle to another agent's sequencer, you need a
typed handle. The macro \chFiveM{uvm\_declare\_p\_sequencer(type)} declares
\code{p\_sequencer} and casts \code{m\_sequencer} into it in \code{pre\_body}:

\begin{uvmcode}
class cordic_cfg_seq extends uvm_sequence #(cordic_item);
  `uvm_object_utils(cordic_cfg_seq)
  `uvm_declare_p_sequencer(cordic_sequencer)
  // ...
  task body();
    // p_sequencer is a cordic_sequencer, not a uvm_sequencer_base
    if (p_sequencer.cfg.stall_pct > 50) ...
  endtask
endclass
\end{uvmcode}

\subsection{Sequence libraries and virtual sequences}

\index{UVM!virtual sequence}%
Two constructions appear in every large testbench and are worth recognising
even though you will not write one this week.

A \emph{sequence library} (\code{uvm\_sequence\_library \#(REQ)}) is a
sequence that picks other sequences and runs them, in random order, in a
random selection, or round-robin. It is how a long soak test is built out of
a dozen small directed sequences.

A \emph{virtual sequence} is a sequence that runs on no protocol at all. It
coordinates several real sequences on several sequencers, which is how you
say ``start the register configuration on the APB agent, wait for it to
finish, then start data traffic on both the input and the output agents at
once''. Its sequencer, a \emph{virtual sequencer}, holds handles to the real
sequencers and drives nothing itself. Virtual sequences are where
system-level tests actually live.

\subsection{The sequencer--driver handshake}
\label{sec:uvm:handshake}

\index{UVM!seq\_item\_port}%
This is the part students get wrong most often, so read
\cref{fig:uvm:handshake} slowly.

\begin{figure}[htbp]
  \centering
  \begin{tikzpicture}[
      font=\sffamily\scriptsize,
      hd/.style={draw=codeframe,fill=svbg,line width=0.8pt,rounded corners=2pt,
                 minimum width=24mm,minimum height=6mm,align=center},
      ll/.style={codeframe,line width=0.8pt,dashed},
      msg/.style={-{Stealth[length=2mm]},codekey,line width=0.8pt},
      ret/.style={-{Stealth[length=2mm]},accentalt,line width=0.8pt,dashed},
      note/.style={font=\sffamily\tiny,text=black!65,align=center}
    ]
    \node[hd] (s) at (0,0)   {sequence\\\texttt{body()}};
    \node[hd] (q) at (4.6,0) {sequencer};
    \node[hd] (d) at (9.2,0) {driver\\\texttt{run\_phase}};

    \draw[ll] (s.south) -- (0,-7.4);
    \draw[ll] (q.south) -- (4.6,-7.4);
    \draw[ll] (d.south) -- (9.2,-7.4);

    \draw[msg] (0,-0.9)  -- node[above,note]{\texttt{start\_item(req)}} (4.6,-0.9);
    \draw[ret] (4.6,-1.6) -- node[above,note]{grant: req is now the active item} (0,-1.6);
    \node[note,anchor=west] at (0.15,-2.25) {\texttt{req.randomize()}};
    \draw[msg] (0,-2.9)  -- node[above,note]{\texttt{finish\_item(req)}} (4.6,-2.9);

    \draw[msg] (9.2,-3.6) -- node[above,note]{\texttt{seq\_item\_port.get\_next\_item(req)}} (4.6,-3.6);
    \draw[ret] (4.6,-4.3) -- node[above,note]{returns the item} (9.2,-4.3);
    \node[note,anchor=east] at (9.05,-4.95) {drive the pins};
    \draw[msg] (9.2,-5.6) -- node[above,note]{\texttt{seq\_item\_port.item\_done()}} (4.6,-5.6);
    \draw[ret] (4.6,-6.3) -- node[above,note]{\texttt{finish\_item} returns; \texttt{body()} continues} (0,-6.3);

    \node[note,anchor=west,text=accentalt] at (-2.6,-7.0)
         {every arrow is a blocking call: nothing here is a signal};
  \end{tikzpicture}
  \caption{One round of the sequencer--driver handshake. The sequence blocks
    in \texttt{finish\_item} from the moment it hands the item over until the
    driver calls \texttt{item\_done}, which is what makes the driver's
    back-pressure propagate all the way back to the stimulus.}
  \label{fig:uvm:handshake}
\end{figure}

On the driver side the loop is always the same three lines:

\begin{uvmcode}[caption={The driver's side of the handshake},label={lst:uvm:drv-loop}]
task run_phase(uvm_phase phase);
  forever begin
    seq_item_port.get_next_item(req);   // blocks until an item arrives
    drive(req);                         // your pin wiggling, takes time
    seq_item_port.item_done();          // releases the sequence
  end
endtask
\end{uvmcode}

\code{req} is not a variable you declare: it is inherited from
\code{uvm\_driver \#(REQ, RSP)}, together with \code{rsp}. The port
\code{seq\_item\_port} is inherited too. All you provide is \code{drive()}.

Notice what this buys you compared with the mailbox in \cref{ch:tb}. There,
\code{req\_mbx.get(t)} removed the transaction from the mailbox and the
generator carried on immediately: the generator ran ahead of the driver,
bounded only by the mailbox depth. Here, the sequence is blocked inside
\code{finish\_item} until \code{item\_done()}, so the stimulus generation
rate is exactly the rate at which the \DUT{} accepts requests. Back-pressure
propagates from the pins to the sequence for free.

\begin{gotcha}{Forgetting \code{item\_done()}}
\begin{uvmcode}
forever begin
  seq_item_port.get_next_item(req);
  drive(req);
  // item_done() missing
end
\end{uvmcode}

The first item is driven correctly. The second \code{get\_next\_item} then
blocks forever, because the sequencer still considers the first item
outstanding. Your simulation drives exactly one transaction and then hangs
until the global timeout, with no error message and a waveform that shows one
perfect handshake followed by silence.

The mirror-image bug is calling \code{item\_done()} twice, or calling it
without a matching \code{get\_next\_item}, which produces a
\code{SEQREQZMB} error from the sequencer (``Item done() called with no
outstanding requests'').

If your driver uses \code{try\_next\_item} instead, remember that
\code{item\_done()} must be called only when \code{try\_next\_item} actually
returned a non-null handle.
\end{gotcha}

\subsection{The pipelined alternative}
\label{sec:uvm:pipelined}

\code{get\_next\_item} is blocking, which is wrong for a driver that must
keep the bus busy every cycle whether or not stimulus is available. The
non-blocking form returns \kw{null} instead of waiting:

\begin{uvmcode}[caption={A driver that never stalls the bus},label={lst:uvm:try-next}]
task run_phase(uvm_phase phase);
  forever begin
    seq_item_port.try_next_item(req);
    if (req == null) begin
      drive_idle();          // hold valid low for one cycle
      @(vif.drv_cb);
    end else begin
      drive(req);
      seq_item_port.item_done();
    end
  end
endtask
\end{uvmcode}

For a genuinely pipelined protocol, where the response to a request arrives
many cycles later and the driver must not block in the meantime, the pattern
is two threads: one calls \code{get\_next\_item}/\code{item\_done()} as fast
as the request channel allows, the other collects responses and returns them
with \code{seq\_item\_port.put\_response(rsp)} or the combined
\code{item\_done(rsp)}. The sequence then reads them with
\code{get\_response(rsp)}. The CORDIC does not need this, because its monitor
observes the responses independently, but you will meet it on any bus
agent.

\begin{ruleofthumb}
If the sequence never needs to see the result, do not use the response path.
Let the monitor observe the result and send it to the scoreboard. Responses
are for sequences that must make a decision based on what came back.
\end{ruleofthumb}

\section{The agent}
\label{sec:uvm:agent}

\index{UVM!agent}%
An agent packages everything that knows about one protocol: a sequencer, a
driver and a monitor, plus whatever configuration they share. It is the unit
of reuse, and it is the one box in \cref{fig:uvm:hierarchy} that has no
counterpart in \cref{ch:tb}.

\code{uvm\_agent} adds exactly one thing to \code{uvm\_component}: a field
called \code{is\_active}, of an enumerated type whose two values are
\code{UVM\_ACTIVE} and \code{UVM\_PASSIVE}. The base class reads it from
the configuration database in its own \code{build\_phase}, and
\code{get\_is\_active()} returns it. An active agent builds all three sub-components and
drives the pins. A passive agent builds only the monitor.

\begin{uvmcode}[caption={The active/passive switch},label={lst:uvm:agent-active}]
function void build_phase(uvm_phase phase);
  super.build_phase(phase);          // reads is_active from the config db
  ap  = new("ap", this);
  mon = cordic_monitor::type_id::create("mon", this);
  if (get_is_active() == UVM_ACTIVE) begin
    sqr = uvm_sequencer #(cordic_item)::type_id::create("sqr", this);
    drv = cordic_driver::type_id::create("drv", this);
  end
endfunction

function void connect_phase(uvm_phase phase);
  mon.ap.connect(ap);
  if (get_is_active() == UVM_ACTIVE)
    drv.seq_item_port.connect(sqr.seq_item_export);
endfunction
\end{uvmcode}

Why does this matter enough to be built into the standard? Because of what
happens when the block is integrated. At block level, your CORDIC agent is
active: it generates the requests. At subsystem level, the CORDIC is fed by
another block, and stimulus now arrives from the design, not from you. You
still want the monitor: you still want the scoreboard to check every
rotation, and you still want the coverage to be collected. So you instantiate
\emph{the same agent}, set it passive, and everything downstream of the
monitor works unchanged.

\begin{uvmcode}
// in the subsystem test's build_phase, before super.build_phase()
uvm_config_db #(uvm_active_passive_enum)::set(
    this, "env.agt", "is_active", UVM_PASSIVE);
\end{uvmcode}

\index{VIP}%
An agent written so that it can be dropped into somebody else's testbench,
self-contained, configured only through the database, with no reference to a
particular \DUT{}, is what the industry calls a \emph{VIP}, a verification
IP. Commercial VIP for AXI, PCIe, USB or DDR is exactly this: a package of
UVM classes with a documented configuration object, a documented sequence
API, and a monitor whose analysis port you connect to your own scoreboard.
Buying VIP instead of writing it is the main commercial reason UVM exists.

\begin{tipbox}{Put the analysis port on the agent, not on the monitor}
In \cref{lst:uvm:agent-active} the agent has its own \code{ap}, and the
monitor's port is connected to it. The environment then connects
\code{agt.ap}, never \code{agt.mon.ap}. This costs one line and it means the
agent's internal structure is private: you can later split the monitor in
two, or add a filter between the monitor and the port, without changing any
environment that uses the agent.
\end{tipbox}

\section{TLM and the analysis path}
\label{sec:uvm:tlm}

\index{UVM!TLM}\index{UVM!analysis port}%
TLM, transaction-level modelling, is the mechanism by which UVM components
exchange objects. You have already met the two forms it takes: the blocking
\code{seq\_item\_port} between sequencer and driver, and the broadcast
analysis port out of the monitor. It is the analysis path that replaces the
mailboxes of \cref{ch:tb}, and it is worth being precise about the
difference.

A \code{mailbox \#(T)} is a point-to-point queue. It has one producer and one
consumer, the producer may block if it is bounded, and the consumer blocks
until something arrives. In \cref{ch:tb} the monitor had to hold two
mailboxes, \code{obs\_mbx} and \code{cov\_mbx}, and put each transaction into
both. To add a third consumer you had to edit the monitor.

A \code{uvm\_analysis\_port \#(T)} is a broadcast. The producer calls
\code{write(t)}, which is a \emph{function}: it never blocks, never consumes
time, and returns immediately. The port forwards the call to every subscriber
connected to it, in an unspecified order. Zero subscribers is legal and does
nothing.

\begin{uvmcode}[caption={Producing on an analysis port},label={lst:uvm:ap}]
class cordic_monitor extends uvm_monitor;
  uvm_analysis_port #(cordic_item) ap;

  function new(string name, uvm_component parent);
    super.new(name, parent);
    ap = new("ap", this);    // ports are NOT factory-created
  endfunction
  // ... in run_phase, when a transaction is complete:
  //     ap.write(it);
endclass
\end{uvmcode}

\subsection{Three ways to consume}

\index{UVM!uvm\_subscriber}%
\textbf{\code{uvm\_subscriber \#(T)}} is the simplest. It is a component with
a built-in \code{analysis\_export} and one pure virtual method,
\code{write(T t)}, which you implement. There is no thread and no queue: your
\code{write} runs inside the monitor's call stack. That makes it ideal for
coverage sampling, which is a function anyway, and unsuitable for anything
that needs to consume time.

\index{UVM!uvm\_analysis\_imp}%
\textbf{\code{uvm\_analysis\_imp \#(T, PARENT)}} is the general form: it
forwards \code{write()} to a method named \code{write} in the parent
component. A component can have only one such method, so a scoreboard that
must listen to two monitors, one on the input and one on the output, cannot
use two plain imps. The macro
\chFiveM{uvm\_analysis\_imp\_decl} solves it by generating a suffixed variant:

\begin{uvmcode}[caption={Two analysis inputs on one scoreboard},label={lst:uvm:imp-decl}]
// At package scope, outside any class:
`uvm_analysis_imp_decl(_req)
`uvm_analysis_imp_decl(_rsp)

class two_port_sb extends uvm_scoreboard;
  `uvm_component_utils(two_port_sb)

  uvm_analysis_imp_req #(cordic_item, two_port_sb) req_imp;
  uvm_analysis_imp_rsp #(cordic_item, two_port_sb) rsp_imp;

  function new(string name, uvm_component parent);
    super.new(name, parent);
    req_imp = new("req_imp", this);
    rsp_imp = new("rsp_imp", this);
  endfunction

  function void write_req(cordic_item t);  /* ... */ endfunction
  function void write_rsp(cordic_item t);  /* ... */ endfunction
endclass
\end{uvmcode}

The macro invents the class \code{uvm\_analysis\_imp\_req} and makes it call
\code{write\_req} instead of \code{write}. It must appear at package or
compilation-unit scope, exactly once per suffix, or you get a duplicate
definition error when two files declare the same suffix.

\index{UVM!uvm\_tlm\_analysis\_fifo}%
\textbf{\code{uvm\_tlm\_analysis\_fifo \#(T)}} is what a scoreboard usually
wants. It is an unbounded FIFO with an \code{analysis\_export} on the input
side and a blocking \code{get()} on the output side. The monitor's
\code{write()} still returns immediately, because the FIFO's \code{write}
only pushes, but the checking code becomes an ordinary \code{run\_phase}
loop:

\begin{uvmcode}
uvm_tlm_analysis_fifo #(cordic_item) fifo;
// ...
task run_phase(uvm_phase phase);
  cordic_item it;
  forever begin
    fifo.get(it);        // blocks until the monitor writes
    check(it);           // may take time, may call other tasks
  end
endtask
\end{uvmcode}

\begin{ruleofthumb}
Coverage collectors get a \code{uvm\_subscriber}. Scoreboards get a
\code{uvm\_tlm\_analysis\_fifo}. Use a bare \code{uvm\_analysis\_imp} only
when you need several differently-named inputs and none of them blocks.
\end{ruleofthumb}

Why does the FIFO matter? Because \code{write()} is a function and cannot
consume time. Any checking that needs to wait, for a matching transaction
from another interface, for a register read, or for a model that itself takes
time, cannot be written inside \code{write()}. Put the item in a FIFO and
do the work in a task. It also decouples the two: if the checker is slow, the
monitor is not held up, because the FIFO is unbounded.

\begin{notebox}{The rest of TLM}
UVM also provides the full TLM-1 port set (\code{uvm\_blocking\_put\_port},
\code{uvm\_nonblocking\_get\_port}, \code{uvm\_transport\_port} and their
exports and imps) and a TLM-2.0 subset with generic payloads and sockets, for
connecting to SystemC models. You will not need them for a block-level
testbench. The three constructions above cover nearly every real use.
\end{notebox}

\section{Reporting}
\label{sec:uvm:reporting}

Stop using \code{\$display}. UVM has four message macros, and using them is
not a style preference: the simulation's pass/fail verdict is computed from
them.

\begin{uvmcode}
`uvm_info(ID, MESSAGE, VERBOSITY)   // informational; filtered by verbosity
`uvm_warning(ID, MESSAGE)           // counted, does not fail the test
`uvm_error(ID, MESSAGE)             // counted, FAILS the test, run continues
`uvm_fatal(ID, MESSAGE)             // prints and terminates immediately
\end{uvmcode}

\code{ID} is a short string used for filtering; \code{get\_type\_name()} is a
sensible default. \code{MESSAGE} is a string, so build it with
\code{\$sformatf}. Note that \chFiveM{uvm\_error} takes no verbosity: errors
are never filtered.

\begin{gotcha}{The macros take no semicolon and no trailing comma}
\begin{uvmcode}
`uvm_info("DRV", $sformatf("n=%0d", n), UVM_LOW)     // correct
`uvm_info("DRV", $sformatf("n=%0d", n), UVM_LOW);    // stray empty statement
`uvm_error("SB", "mismatch", UVM_LOW)                // wrong arity: compile error
\end{uvmcode}
The stray semicolon is harmless after a statement but becomes a syntax error
if the macro is the only thing inside an \kw{if} without \kw{begin}/\kw{end},
which is precisely where you will hit it. Get into the habit of writing the
macros without a semicolon, and of putting \kw{begin}/\kw{end} around any
\kw{if} body that contains one.
\end{gotcha}

\subsection{Verbosity}

\index{UVM!verbosity}%
Each \chFiveM{uvm\_info} carries a verbosity level:
\code{UVM\_NONE} (0), \code{UVM\_LOW} (100), \code{UVM\_MEDIUM} (200),
\code{UVM\_HIGH} (300), \code{UVM\_FULL} (400), \code{UVM\_DEBUG} (500). A
message is printed if its level is \emph{less than or equal to} the current
threshold, which defaults to \code{UVM\_MEDIUM}. So \code{UVM\_NONE} always
prints and \code{UVM\_DEBUG} almost never does.

The threshold is set on the command line,

\begin{shcode}
vsim -c +UVM_VERBOSITY=UVM_HIGH ...
\end{shcode}

or, more usefully, per subtree at run time:

\begin{uvmcode}
// only the driver becomes chatty; everything else stays at MEDIUM
env.agt.drv.set_report_verbosity_level_hier(UVM_HIGH);

// or by message ID, anywhere in the tree
uvm_top.set_report_id_verbosity_hier("HANDSHAKE", UVM_FULL);
\end{uvmcode}

\begin{tipbox}{Choose verbosity levels with a policy}
A convention that works: \code{UVM\_NONE} for the end-of-test summary that
must always appear; \code{UVM\_LOW} for once-per-test events (test started,
configuration chosen, sequence finished); \code{UVM\_MEDIUM} for
once-per-transaction; \code{UVM\_HIGH} for once-per-clock-cycle detail. A
regression then runs at the default and produces a readable log, and a
failing test can be re-run at \code{UVM\_HIGH} for the same seed without
recompiling.
\end{tipbox}

\subsection{The verdict}

\index{UVM!pass/fail verdict}%
At the end of the run, \code{uvm\_report\_server} prints a summary:

\begin{txtcode}
--- UVM Report Summary ---

** Report counts by severity
UVM_INFO :  427
UVM_WARNING :  0
UVM_ERROR :    0
UVM_FATAL :    0
** Report counts by id
[RNTST]            1
[cordic_driver]  500
[cordic_scoreboard] 2
[TEST_DONE]        1
\end{txtcode}

Read the \code{UVM\_ERROR} and \code{UVM\_FATAL} lines: those two numbers are
the verdict. A regression script greps for exactly this block. Note what this
implies:

\begin{itemize}
  \item A test that calls \code{\$finish} directly bypasses the summary and
        produces no verdict. Never call \code{\$finish} in a UVM testbench.
  \item A check written with \code{\$error} instead of \chFiveM{uvm\_error}
        is not counted, so a failing test reports zero errors and the
        regression records a pass.
  \item A test in which nothing happened also reports zero errors. This is
        why the scoreboard in \cref{sec:uvm:complete} raises an error when
        it has checked nothing: ``no failures'' and ``no checks'' must not look
        the same.
\end{itemize}

\index{UVM!UVM\_TESTNAME}%
Two command-line switches complete the picture. The switch
\code{+UVM\_TESTNAME} names the \code{uvm\_test} subclass that
\code{run\_test()} builds, so one compiled image runs every test in the
suite. And
\code{+UVM\_MAX\_QUIT\_COUNT} stops the run after a given number of errors,
which keeps a broken test from producing a gigabyte of log.

\section{A complete UVM testbench for the CORDIC}
\label{sec:uvm:complete}

Everything so far, assembled. This is a genuinely runnable testbench: type it
in, compile it against \file{cordic\_pkg.sv}, \file{cordic\_if.sv} and
\file{cordic\_rot.sv} from \cref{ch:tb}, and it will find bugs in the CORDIC.
Each listing names the class from \file{tb\_cordic\_pkg.sv} that it replaces.

Everything except the top-level module lives in one package:

\begin{uvmcode}[caption={Package header},label={lst:uvm:pkg}]
`timescale 1ns/1ps

package cordic_uvm_pkg;

  import uvm_pkg::*;
  `include "uvm_macros.svh"
  import cordic_pkg::*;

  // The CORDIC processing gain and the comparison tolerance, exactly as in
  // the hand-written testbench of the previous chapter.
  localparam real K_GAIN  = 1.6467602581210656;
  localparam int  TOL_LSB = 12;
\end{uvmcode}

\subsection{The sequence item}

Replaces \code{cordic\_txn}. Compare it with the hand-written class of
\cref{lst:tb:txn}: the constraints are identical, in the ordinary \SV{} form
of \cref{ch:toolkit}, the \code{convert2string} is identical, and everything
else is base-class machinery.

\begin{uvmcode}[caption={\code{cordic\_item}: replaces \code{cordic\_txn}},label={lst:uvm:item}]
  class cordic_item extends uvm_sequence_item;

    rand data_t  x0;
    rand data_t  y0;
    rand angle_t theta;

    data_t x;              // observed, filled in by the monitor
    data_t y;

    `uvm_object_utils_begin(cordic_item)
      `uvm_field_int(x0,    UVM_DEFAULT | UVM_DEC)
      `uvm_field_int(y0,    UVM_DEFAULT | UVM_DEC)
      `uvm_field_int(theta, UVM_DEFAULT | UVM_DEC)
      `uvm_field_int(x,     UVM_DEFAULT | UVM_DEC)
      `uvm_field_int(y,     UVM_DEFAULT | UVM_DEC)
    `uvm_object_utils_end

    constraint c_theta {
      theta inside {[-PI_EXT : PI_EXT]};
    }
    constraint c_magnitude {
      x0 inside {[-6000 : 6000]};
      y0 inside {[-6000 : 6000]};
    }
    constraint c_corners {
      theta dist {
        0                :/ 1,
        PI_EXT           :/ 1,
        -PI_EXT          :/ 1,
        PI_EXT / 2       :/ 1,
        -PI_EXT / 2      :/ 1,
        [-PI_EXT:PI_EXT] :/ 20
      };
    }

    function new(string name = "cordic_item");
      super.new(name);
    endfunction

    function string convert2string();
      return $sformatf(
        "x0=%0d (%.4f) y0=%0d (%.4f) theta=%0d (%.4f rad) -> x=%0d y=%0d",
        x0, q_to_real(x0, QF), y0, q_to_real(y0, QF),
        theta, q_to_real(theta, QA), x, y);
    endfunction

  endclass : cordic_item
\end{uvmcode}

\subsection{The sequence}

Replaces \code{cordic\_env::send\_random()}. Note that it is now a class, so
a directed test can extend it, and a regression can select it by name.

\begin{uvmcode}[caption={\code{cordic\_random\_seq} and a directed sequence},label={lst:uvm:seq}]
  class cordic_random_seq extends uvm_sequence #(cordic_item);
    `uvm_object_utils(cordic_random_seq)

    rand int unsigned n_items = 200;
    constraint c_n { n_items inside {[1 : 5000]}; }

    function new(string name = "cordic_random_seq");
      super.new(name);
    endfunction

    task body();
      cordic_item it;
      `uvm_info(get_type_name(),
                $sformatf("sending %0d random items", n_items), UVM_LOW)
      repeat (n_items) begin
        it = cordic_item::type_id::create("it");
        start_item(it);
        if (!it.randomize())
          `uvm_fatal(get_type_name(), "item randomization failed")
        finish_item(it);
      end
    endtask

  endclass : cordic_random_seq

  // The directed counterpart: the corner angles, with a unit-length seed so
  // that the output is (cos theta, sin theta) directly.
  class cordic_corner_seq extends uvm_sequence #(cordic_item);
    `uvm_object_utils(cordic_corner_seq)

    function new(string name = "cordic_corner_seq");
      super.new(name);
    endfunction

    task body();
      cordic_item it;
      angle_t angles[$] = '{0, PI_EXT, -PI_EXT, PI_EXT/2, -PI_EXT/2,
                            PI_EXT/4, -PI_EXT/4};
      foreach (angles[i]) begin
        it = cordic_item::type_id::create("it");
        start_item(it);
        it.x0    = data_t'(X0_UNIT);
        it.y0    = '0;
        it.theta = angles[i];
        finish_item(it);
      end
    endtask

  endclass : cordic_corner_seq
\end{uvmcode}

\subsection{The driver}

Replaces \code{cordic\_drv}. The \code{send()} and \code{reset()} bodies are
line-for-line the ones from \cref{ch:tb}; what changed is where the
transactions come from and who starts the thread.

\begin{uvmcode}[caption={\code{cordic\_driver}: replaces \code{cordic\_drv}},label={lst:uvm:drv}]
  class cordic_driver extends uvm_driver #(cordic_item);
    `uvm_component_utils(cordic_driver)

    virtual cordic_if vif;
    int unsigned      stall_pct = 25;   // random back-pressure on rsp_ready

    function new(string name, uvm_component parent);
      super.new(name, parent);
    endfunction

    function void build_phase(uvm_phase phase);
      super.build_phase(phase);
      if (!uvm_config_db #(virtual cordic_if)::get(this, "", "vif", vif))
        `uvm_fatal("NOVIF",
                   {"no virtual interface for ", get_full_name(), ".vif"})
      void'(uvm_config_db #(int unsigned)::get(this, "", "stall_pct",
                                               stall_pct));
    endfunction

    task run_phase(uvm_phase phase);
      do_reset();
      fork
        response_backpressure();
      join_none
      forever begin
        seq_item_port.get_next_item(req);
        drive(req);
        seq_item_port.item_done();
      end
    endtask

    protected task do_reset();
      vif.rst_n            <= 1'b0;
      vif.drv_cb.req_valid <= 1'b0;
      vif.drv_cb.req_x     <= '0;
      vif.drv_cb.req_y     <= '0;
      vif.drv_cb.req_theta <= '0;
      vif.drv_cb.rsp_ready <= 1'b0;
      repeat (5) @(vif.drv_cb);
      vif.rst_n <= 1'b1;
      @(vif.drv_cb);
    endtask

    // Drive one request and hold it until the DUT accepts it.  Marked
    // `virtual' so that a factory override can replace it -- this is what
    // the error-injecting driver shown earlier extends.
    protected virtual task drive(cordic_item it);
      vif.drv_cb.req_valid <= 1'b1;
      vif.drv_cb.req_x     <= it.x0;
      vif.drv_cb.req_y     <= it.y0;
      vif.drv_cb.req_theta <= it.theta;
      do @(vif.drv_cb); while (vif.drv_cb.req_ready !== 1'b1);
      vif.drv_cb.req_valid <= 1'b0;
      `uvm_info(get_type_name(),
                {"drove ", it.convert2string()}, UVM_HIGH)
    endtask

    protected task response_backpressure();
      forever begin
        if ($urandom_range(99) < stall_pct) begin
          vif.drv_cb.rsp_ready <= 1'b0;
          repeat ($urandom_range(3) + 1) @(vif.drv_cb);
        end
        vif.drv_cb.rsp_ready <= 1'b1;
        @(vif.drv_cb);
      end
    endtask

  endclass : cordic_driver
\end{uvmcode}

\subsection{The monitor}

Replaces \code{cordic\_mon}. The two mailboxes have become one analysis port.

\begin{uvmcode}[caption={\code{cordic\_monitor}: replaces \code{cordic\_mon}},label={lst:uvm:mon}]
  class cordic_monitor extends uvm_monitor;
    `uvm_component_utils(cordic_monitor)

    virtual cordic_if                vif;
    uvm_analysis_port #(cordic_item) ap;

    protected cordic_item pending[$];   // a queue: accepted, not answered

    function new(string name, uvm_component parent);
      super.new(name, parent);
      ap = new("ap", this);
    endfunction

    function void build_phase(uvm_phase phase);
      super.build_phase(phase);
      if (!uvm_config_db #(virtual cordic_if)::get(this, "", "vif", vif))
        `uvm_fatal("NOVIF",
                   {"no virtual interface for ", get_full_name(), ".vif"})
    endfunction

    task run_phase(uvm_phase phase);
      cordic_item it;
      forever begin
        @(vif.mon_cb);
        if (vif.rst_n !== 1'b1) begin
          pending.delete();
          continue;
        end

        // Request accepted: remember what went in.
        if (vif.mon_cb.req_valid === 1'b1 &&
            vif.mon_cb.req_ready === 1'b1) begin
          it = cordic_item::type_id::create("it");
          it.x0    = vif.mon_cb.req_x;
          it.y0    = vif.mon_cb.req_y;
          it.theta = vif.mon_cb.req_theta;
          pending.push_back(it);
        end

        // Response accepted: complete the oldest outstanding request and
        // broadcast it.  Every subscriber sees the same handle, so nobody
        // downstream may modify it.
        if (vif.mon_cb.rsp_valid === 1'b1 &&
            vif.mon_cb.rsp_ready === 1'b1) begin
          if (pending.size() == 0) begin
            `uvm_error(get_type_name(),
                       "response with no outstanding request")
          end else begin
            it   = pending.pop_front();
            it.x = vif.mon_cb.rsp_x;
            it.y = vif.mon_cb.rsp_y;
            ap.write(it);
          end
        end
      end
    endtask

    // Runs after the run phase: anything still pending is a lost response.
    function void check_phase(uvm_phase phase);
      super.check_phase(phase);
      if (pending.size() != 0)
        `uvm_error(get_type_name(),
                   $sformatf("%0d requests never produced a response",
                             pending.size()))
    endfunction

  endclass : cordic_monitor
\end{uvmcode}

The \code{check\_phase} is new. In \cref{ch:tb} a lost response simply meant
the drain loop timed out and \code{\$fatal} fired with an unhelpful message.
Here the check is explicit and it names the number of lost transactions.

\subsection{The scoreboard}

Replaces \code{cordic\_sb}. The reference model is unchanged: it is written
against the specification, in real arithmetic, and knows nothing about
micro-rotations.

\begin{uvmcode}[caption={\code{cordic\_scoreboard}: replaces \code{cordic\_sb}},label={lst:uvm:sb}]
  class cordic_scoreboard extends uvm_scoreboard;
    `uvm_component_utils(cordic_scoreboard)

    uvm_tlm_analysis_fifo #(cordic_item) fifo;

    int unsigned n_checked = 0;
    int unsigned n_failed  = 0;
    real         worst_err = 0.0;

    function new(string name, uvm_component parent);
      super.new(name, parent);
      fifo = new("fifo", this);
    endfunction

    // The golden model: the specification, not the implementation.
    function void predict(input cordic_item it,
                          output real x_ref, output real y_ref);
      real th = q_to_real(it.theta, QA);
      real xr = q_to_real(it.x0,    QF);
      real yr = q_to_real(it.y0,    QF);
      x_ref = K_GAIN * (xr * $cos(th) - yr * $sin(th));
      y_ref = K_GAIN * (xr * $sin(th) + yr * $cos(th));
    endfunction

    task run_phase(uvm_phase phase);
      cordic_item it;
      real x_ref, y_ref, ex, ey;
      forever begin
        fifo.get(it);
        predict(it, x_ref, y_ref);

        ex = (q_to_real(it.x, QF) - x_ref) * (2.0 ** QF);   // error in LSB
        ey = (q_to_real(it.y, QF) - y_ref) * (2.0 ** QF);
        if (ex < 0.0) ex = -ex;
        if (ey < 0.0) ey = -ey;
        if (ex > worst_err) worst_err = ex;
        if (ey > worst_err) worst_err = ey;

        n_checked++;
        if (ex > real'(TOL_LSB) || ey > real'(TOL_LSB)) begin
          n_failed++;
          `uvm_error(get_type_name(),
            $sformatf("MISMATCH %s want x=%.1f y=%.1f err %.2f / %.2f LSB",
                      it.convert2string(),
                      x_ref * (2.0 ** QF), y_ref * (2.0 ** QF), ex, ey))
        end
      end
    endtask

    function void check_phase(uvm_phase phase);
      super.check_phase(phase);
      // "no failures" and "no checks" must not look the same.
      if (n_checked == 0)
        `uvm_error(get_type_name(), "no transactions were checked")
      if (!fifo.is_empty())
        `uvm_error(get_type_name(), "transactions left unchecked in the fifo")
    endfunction

    function void report_phase(uvm_phase phase);
      `uvm_info(get_type_name(),
        $sformatf("%0d checked, %0d failed, worst error %.2f LSB (tol %0d)",
                  n_checked, n_failed, worst_err, TOL_LSB), UVM_NONE)
    endfunction

  endclass : cordic_scoreboard
\end{uvmcode}

\subsection{The coverage collector}

Replaces \code{cordic\_cov}. The covergroup is character-for-character the
one from \cref{ch:tb}; what disappeared is the mailbox and the
\code{forever} loop, because \code{write()} is called for us.

\begin{uvmcode}[caption={\code{cordic\_coverage}: replaces \code{cordic\_cov}},label={lst:uvm:cov}]
  class cordic_coverage extends uvm_subscriber #(cordic_item);
    `uvm_component_utils(cordic_coverage)

    angle_t sampled_theta;
    data_t  sampled_x0;

    covergroup cg_stimulus;
      option.per_instance = 1;

      cp_quadrant : coverpoint sampled_theta {
        bins q3_neg = {[-PI_EXT   : -PI_EXT/2 - 1]};
        bins q4_neg = {[-PI_EXT/2 : -1]};
        bins zero   = {0};
        bins q1_pos = {[1         :  PI_EXT/2 - 1]};
        bins q2_pos = {[PI_EXT/2  :  PI_EXT]};
      }
      cp_sign_x : coverpoint sampled_x0 {
        bins neg  = {[-32768 : -1]};
        bins zero = {0};
        bins pos  = {[1 : 32767]};
      }
      x_quadrant_sign : cross cp_quadrant, cp_sign_x;
    endgroup

    function new(string name, uvm_component parent);
      super.new(name, parent);
      cg_stimulus = new();     // covergroups are not factory objects
    endfunction

    // uvm_subscriber declares this pure virtual; the argument must be
    // named `t'.  It is a function: it cannot consume time.
    function void write(cordic_item t);
      sampled_theta = t.theta;
      sampled_x0    = t.x0;
      cg_stimulus.sample();
    endfunction

    function void report_phase(uvm_phase phase);
      `uvm_info(get_type_name(),
                $sformatf("functional coverage %.1f %%",
                          cg_stimulus.get_coverage()), UVM_NONE)
    endfunction

  endclass : cordic_coverage
\end{uvmcode}

\subsection{The agent and the environment}

Replaces \code{cordic\_env}, split into the reusable half (the agent) and the
\DUT{}-specific half (the environment).

\begin{uvmcode}[caption={\code{cordic\_agent} and \code{cordic\_env}},label={lst:uvm:env}]
  class cordic_agent extends uvm_agent;
    `uvm_component_utils(cordic_agent)

    cordic_driver                    drv;
    uvm_sequencer #(cordic_item)     sqr;
    cordic_monitor                   mon;
    uvm_analysis_port #(cordic_item) ap;

    function new(string name, uvm_component parent);
      super.new(name, parent);
    endfunction

    function void build_phase(uvm_phase phase);
      super.build_phase(phase);      // reads is_active from the config db
      ap  = new("ap", this);
      mon = cordic_monitor::type_id::create("mon", this);
      if (get_is_active() == UVM_ACTIVE) begin
        sqr = uvm_sequencer #(cordic_item)::type_id::create("sqr", this);
        drv = cordic_driver::type_id::create("drv", this);
      end
    endfunction

    function void connect_phase(uvm_phase phase);
      mon.ap.connect(ap);            // republish through the agent's port
      if (get_is_active() == UVM_ACTIVE)
        drv.seq_item_port.connect(sqr.seq_item_export);
    endfunction

  endclass : cordic_agent


  class cordic_env extends uvm_env;
    `uvm_component_utils(cordic_env)

    cordic_agent      agt;
    cordic_scoreboard sb;
    cordic_coverage   cov;

    function new(string name, uvm_component parent);
      super.new(name, parent);
    endfunction

    function void build_phase(uvm_phase phase);
      super.build_phase(phase);
      agt = cordic_agent::type_id::create("agt", this);
      sb  = cordic_scoreboard::type_id::create("sb",  this);
      cov = cordic_coverage::type_id::create("cov", this);
    endfunction

    // One producer, two consumers.  Adding a third costs one line here and
    // nothing at all in the monitor.
    function void connect_phase(uvm_phase phase);
      agt.ap.connect(sb.fifo.analysis_export);
      agt.ap.connect(cov.analysis_export);
    endfunction

  endclass : cordic_env
\end{uvmcode}

\subsection{The tests}

New: in \cref{ch:tb} the test was the top-level module. Here it is a class,
so the same compiled image runs all of them.

\begin{uvmcode}[caption={A base test and two derived tests},label={lst:uvm:test}]
  class cordic_base_test extends uvm_test;
    `uvm_component_utils(cordic_base_test)

    cordic_env env;

    function new(string name, uvm_component parent);
      super.new(name, parent);
    endfunction

    function void build_phase(uvm_phase phase);
      super.build_phase(phase);
      env = cordic_env::type_id::create("env", this);
    endfunction

    function void end_of_elaboration_phase(uvm_phase phase);
      super.end_of_elaboration_phase(phase);
      uvm_top.print_topology();      // the tree, once, at time zero
    endfunction

    // Common to every derived test: give the DUT pipeline time to empty
    // after the last request, before the run phase is killed.
    protected virtual function void set_drain(uvm_phase phase);
      phase.phase_done.set_drain_time(this, 500ns);
    endfunction

  endclass : cordic_base_test


  class cordic_random_test extends cordic_base_test;
    `uvm_component_utils(cordic_random_test)

    function new(string name, uvm_component parent);
      super.new(name, parent);
    endfunction

    task run_phase(uvm_phase phase);
      cordic_random_seq seq;
      phase.raise_objection(this, "random stimulus");
      set_drain(phase);
      seq = cordic_random_seq::type_id::create("seq");
      if (!seq.randomize() with { n_items == 500; })
        `uvm_fatal(get_type_name(), "sequence randomization failed")
      seq.start(env.agt.sqr);
      phase.drop_objection(this, "random stimulus");
    endtask

  endclass : cordic_random_test


  class cordic_corner_test extends cordic_base_test;
    `uvm_component_utils(cordic_corner_test)

    function new(string name, uvm_component parent);
      super.new(name, parent);
    endfunction

    task run_phase(uvm_phase phase);
      cordic_corner_seq dseq;
      cordic_random_seq rseq;
      phase.raise_objection(this, "corners then random");
      set_drain(phase);
      dseq = cordic_corner_seq::type_id::create("dseq");
      dseq.start(env.agt.sqr);
      rseq = cordic_random_seq::type_id::create("rseq");
      void'(rseq.randomize() with { n_items == 100; });
      rseq.start(env.agt.sqr);
      phase.drop_objection(this, "corners then random");
    endtask

  endclass : cordic_corner_test

endpackage : cordic_uvm_pkg
\end{uvmcode}

\subsection{The top-level module}

The only non-class file. It holds the clock, the interface, the \DUT{}, the
one configuration-database write and the call to \code{run\_test()}.

\begin{uvmcode}[caption={\code{tb\_cordic\_uvm}: the top-level module},label={lst:uvm:top}]
`timescale 1ns/1ps

module tb_cordic_uvm;

  import uvm_pkg::*;
  `include "uvm_macros.svh"
  import cordic_pkg::*;
  import cordic_uvm_pkg::*;

  // ---- clock ------------------------------------------------------------
  logic clk = 1'b0;
  always #5ns clk = ~clk;          // 100 MHz

  // ---- the pins ---------------------------------------------------------
  cordic_if vif (.clk(clk));

  // ---- the DUT ----------------------------------------------------------
  cordic_rot #(.ITER(N)) dut (
    .clk_i       (clk),
    .rst_ni      (vif.rst_n),
    .req_valid_i (vif.req_valid),
    .req_ready_o (vif.req_ready),
    .req_x_i     (vif.req_x),
    .req_y_i     (vif.req_y),
    .req_theta_i (vif.req_theta),
    .rsp_valid_o (vif.rsp_valid),
    .rsp_ready_i (vif.rsp_ready),
    .rsp_x_o     (vif.rsp_x),
    .rsp_y_o     (vif.rsp_y)
  );

  // ---- hand the interface to the class world, then start UVM ------------
  initial begin
    uvm_config_db #(virtual cordic_if)::set(null, "*", "vif", vif);
    uvm_top.set_timeout(1ms, 1);
    run_test();                     // never returns
  end

  // Portable waveform dump, off by default: pass +DUMP to enable it.  Name
  // the scope; $dumpvars(0, tb_cordic_uvm) would dump the whole hierarchy.
  initial if ($test$plusargs("DUMP")) begin
    $dumpfile("tb_cordic_uvm.vcd");
    $dumpvars(1, dut);
  end

endmodule : tb_cordic_uvm
\end{uvmcode}

\begin{notebox}{What \code{run\_test()} actually does}
\code{run\_test()} is a task in \code{uvm\_pkg}. It reads
\code{+UVM\_TESTNAME}, asks the factory for an object of that type, gives it
the name \code{uvm\_test\_top} and no parent, and then runs the phase
schedule over the tree that the test's \code{build\_phase} creates. When the
last phase finishes it prints the report summary and calls \code{\$finish}
itself. It does not return, which is why nothing may follow it in the
\kw{initial} block.

If \code{+UVM\_TESTNAME} names a class that is not registered with the
factory, you get \code{UVM\_FATAL ... [INVTST] Requested test from
command line +UVM\_TESTNAME=... not found}. The usual cause is a missing
\chFiveM{uvm\_component\_utils}, or a package that was compiled but never
imported by the top-level module, so its static registrations never ran.
\end{notebox}

\section{The register layer}
\label{sec:uvm:reg}

\index{UVM!register layer}%
The CORDIC has no registers, so this section is a map rather than a
tutorial. Any block with a control interface has a register map, and UVM
provides a modelling layer for it.

\begin{description}
  \item[\code{uvm\_reg\_field}] one field: width, offset, access policy
        (\code{RW}, \code{RO}, \code{W1C}, \code{RC}, and about twenty
        others), reset value.
  \item[\code{uvm\_reg}] one register: a name, an offset, and its fields.
        Each register carries a \emph{mirror}, the model's belief about the
        hardware's current value, and a \emph{desired} value.
  \item[\code{uvm\_reg\_block}] a group of registers and other blocks, with
        one or more address maps. The block is the unit that matches an
        IP-XACT component.
  \item[\code{uvm\_reg\_adapter}] the translator. It converts a generic
        \code{uvm\_reg\_bus\_op} into your bus agent's sequence item and
        back. You write one adapter per bus protocol, ever.
  \item[\code{uvm\_reg\_predictor}] listens on the bus monitor's analysis
        port and updates the mirror from observed traffic, so that the model
        stays correct even when somebody else writes the registers.
\end{description}

\index{UVM!back-door access}%
With a model in place, a sequence stops talking about addresses:

\begin{uvmcode}
uvm_status_e status;
uvm_reg_data_t value;

// front door: a real bus transaction, through the adapter and the driver
regmodel.ctrl.write(status, 32'h0000_0003, UVM_FRONTDOOR);
regmodel.ctrl.read (status, value,         UVM_FRONTDOOR);

// back door: a direct hierarchical poke, zero simulation time
regmodel.ctrl.write(status, 32'h0000_0003, UVM_BACKDOOR);
\end{uvmcode}

A \emph{front-door} access produces real cycles on the bus, exercises the bus
interface, and takes time. A \emph{back-door} access reaches into the design
hierarchy with a hierarchical reference (the path is supplied by
\code{add\_hdl\_path}) and changes the register directly, in zero time. Use
the front door to verify the register interface, and the back door to set up
a state you do not want to spend a thousand cycles reaching.

What the model buys you: a set of ready-made sequences that check every reset
value, every writable bit and every access policy without a line of test code
(\code{uvm\_reg\_hw\_reset\_seq} and its relatives); a mirror that lets a
scoreboard ask ``what is the current mode?'' without snooping the bus; and
tests written against register names, which survive an address-map change.

\begin{notebox}{Nobody writes a register model by hand}
A medium-sized block has a few hundred registers, and the model is roughly
twenty lines per register. Generating it is standard practice. The input is
whatever the project treats as the specification: an IP-XACT
\code{.xml} component description, a SystemRDL file, or, very commonly, a
spreadsheet exported to CSV. Tools such as \tool{PeakRDL},
\tool{Semifore CSRCompiler}, \tool{IDesignSpec} or a home-made Python script
emit the UVM model, the \RTL{}, the C header and the documentation from the
same source. If you find yourself editing a \code{uvm\_reg\_block} by hand,
somebody has lost the generator.
\end{notebox}

\section{Running it}
\label{sec:uvm:running}

\index{QuestaSim!UVM}%
UVM is a package. Compiling it is compiling one file with one include
directory, and every simulator has a shortcut for it.

\subsection{QuestaSim}

The explicit form, which shows what is happening. It is the flow of
\cref{sec:tb:run} with one extra compile at the front:

\begin{shcode}
vlib work
vmap work work

# 1. the UVM library itself (once; $UVM_HOME points into the Questa install)
vlog -sv +incdir+$UVM_HOME/src $UVM_HOME/src/uvm_pkg.sv

# 2. the design and the testbench
vlog -sv +incdir+$UVM_HOME/src cordic_pkg.sv cordic_if.sv cordic_rot.sv
vlog -sv +incdir+$UVM_HOME/src cordic_uvm_pkg.sv tb_cordic_uvm.sv

# 3. optimise, keeping enough visibility for the factory and the waveform
vopt -o tb_opt +acc=rnb work.tb_cordic_uvm

# 4. run one test
vsim -c -sv_seed random work.tb_opt \
     +UVM_TESTNAME=cordic_random_test \
     +UVM_VERBOSITY=UVM_MEDIUM \
     -do uvm.do
\end{shcode}

The short form uses the built-in library, which is pre-compiled and faster,
and lets \code{vsim} run \code{vopt} for you:

\begin{shcode}
vlog -sv -uvm cordic_pkg.sv cordic_if.sv cordic_rot.sv \
              cordic_uvm_pkg.sv tb_cordic_uvm.sv
vsim -c -uvm -voptargs="+acc=rnb" tb_cordic_uvm \
     +UVM_TESTNAME=cordic_random_test -sv_seed random -do uvm.do
\end{shcode}

The \code{do} file is shorter than the one in \cref{sec:tb:run}, because UVM
computes the verdict itself and \code{run\_test()} calls \code{\$finish} when
the last phase is over:

\begin{tclcode}
# uvm.do -- batch execution of a UVM test
onerror {quit -code 3}

# log the DUT only.  `log -r /*' would also record every UVM class member.
log -r /tb_cordic_uvm/dut/*

run -all
quit -code 0
\end{tclcode}

Note the scope on \code{log -r}. Logging the whole hierarchy in a UVM testbench
records the class world as well as the pins, and a \file{.wlf} that would have
been tens of megabytes becomes hundreds. \tool{QuestaSim} shows that
\file{.wlf} in its own wave window; the \code{\$dumpfile}/\code{\$dumpvars}
pair in \cref{lst:uvm:top} is the portable alternative, writing a VCD that
\tool{GTKWave} and \tool{Surfer} can open, and the same warning applies to its
first argument. \cref{sec:tb:waves} covers the choice.

Use \code{-uvmhome <path>} if you need a specific UVM version rather than the
one bundled with the tool. Mixing a hand-compiled \file{uvm\_pkg.sv} with the
tool's built-in one produces duplicate-definition errors that are tedious to
untangle, so pick one and be consistent across the whole project.

\code{-sv\_seed random} is important. Without it every run uses the same
seed, and a constrained-random testbench that always generates the same
stimulus is a directed testbench with extra steps. Print the seed (Questa
does so in the banner) so that a failing run can be reproduced exactly with
\code{-sv\_seed <n>}.

\subsection{The other two}

\begin{shcode}
# Cadence Xcelium
xrun -uvm -sv cordic_pkg.sv cordic_if.sv cordic_rot.sv \
     cordic_uvm_pkg.sv tb_cordic_uvm.sv \
     +UVM_TESTNAME=cordic_random_test -svseed random

# Synopsys VCS
vcs -sverilog -ntb_opts uvm -full64 cordic_pkg.sv cordic_if.sv \
    cordic_rot.sv cordic_uvm_pkg.sv tb_cordic_uvm.sv -o simv
./simv +UVM_TESTNAME=cordic_random_test +ntb_random_seed_automatic
\end{shcode}

\subsection{Reading the output}

A successful run ends like this:

\begin{txtcode}
UVM_INFO cordic_scoreboard.sv(78) @ 26505 ns:
  uvm_test_top.env.sb [cordic_scoreboard]
  500 checked, 0 failed, worst error 3.41 LSB (tol 12)
UVM_INFO cordic_coverage.sv(52) @ 26505 ns:
  uvm_test_top.env.cov [cordic_coverage]
  functional coverage 93.3 %
UVM_INFO .../uvm_report_server.svh(847) @ 26505 ns:
  reporter [UVM/REPORT/SERVER]
--- UVM Report Summary ---

** Report counts by severity
UVM_INFO :    4
UVM_WARNING : 0
UVM_ERROR :   0
UVM_FATAL :   0
\end{txtcode}

Every message carries the file, the line, the simulation time and the full
hierarchical name of the component that produced it. That last field is the
one you will use: \code{uvm\_test\_top.env.sb} tells you immediately which
instance complained, which in a testbench with six agents is not obvious from
the message text.

\begin{tipbox}{Make the regression script read the right thing}
Grep for the two lines \code{UVM\_ERROR :} and \code{UVM\_FATAL :} in the
report summary and require both counts to be zero. Do not grep for the word
``error'' anywhere in the log: a message containing the word ``error'' in its
text (\code{worst error 3.41 LSB}, above) will fail a passing test. Do not
rely on the simulator's exit status either, because a UVM run that ends with
errors still exits normally.
\end{tipbox}

\section{When not to use UVM}
\label{sec:uvm:when-not}

\index{UVM!when not to use}%
This section exists because everything before it makes UVM look inevitable,
and it is not.

\textbf{A block with three ports and one protocol does not need it.} If the
whole verification effort is ``drive a valid/ready channel, check the output
against a model'', you will spend more lines on UVM infrastructure than on
verification. The testbench in \cref{ch:tb} does that job in about 470 lines that a
new team member reads in twenty minutes.

\textbf{A three-week project does not need it.} The break-even point for UVM
is somewhere around the moment you write your third test that reuses the
second test's environment. If the project ends before that, the investment
never pays back.

\textbf{A course exercise does not need it.} A laboratory assignment is
graded on whether the design works, and a UVM testbench in which the student
does not understand the phase order teaches nothing about verification. Learn
UVM \emph{after} you can write \cref{ch:tb}'s testbench from memory, not
instead.

\textbf{A single-person project rarely needs it.} Much of UVM's value is
social: a shared vocabulary, a shared structure, an agent one engineer can
hand to another. With one engineer, the shared vocabulary is your own habits.

Where UVM does pay:

\begin{itemize}
  \item \textbf{Reuse across blocks.} Five blocks on the same bus share one
        bus agent, one adapter and one register-model generator.
  \item \textbf{Reuse across levels.} The block-level agent goes passive and
        the same monitor, scoreboard and coverage run at subsystem and chip
        level.
  \item \textbf{Reuse across projects.} A protocol agent written once is used
        for a decade.
  \item \textbf{A team larger than one.} Structure that is enforced by a
        library is structure that survives people joining and leaving.
  \item \textbf{Commercial VIP.} If you buy an AXI or PCIe VIP, it is UVM,
        and your testbench must be too.
\end{itemize}

\subsection{The alternatives, fairly}

\index{OSVVM}\index{UVVM}\index{cocotb}\index{pyuvm}%
\begin{description}
  \item[A hand-written layered testbench.] What you built in \cref{ch:tb}.
        Full \SV{} power, meaning classes, constrained random, functional
        coverage and assertions, with none of the library. It scales further than people
        expect, and the point at which it stops scaling is usually the point
        at which you need reuse across teams. Many production testbenches for
        small blocks are exactly this.
  \item[cocotb, with or without pyuvm.] The testbench is Python; the
        simulator is driven through VPI/VHPI. Stimulus, checking and the
        reference model are written in a language with real libraries:
        \code{numpy} for the CORDIC's golden model, \code{pytest} for the
        regression. \code{pyuvm} is a faithful port of the UVM class library
        to Python, so the six boxes and the phases keep their names. It drives
        \tool{QuestaSim} through the FLI and VPI, and it also drives
        \tool{Verilator} and \tool{Icarus Verilog}, so a cocotb testbench can
        run without a licence. \refsimulationbook builds one.
  \item[OSVVM, for \VHDL{}.] A \VHDL{} library providing randomisation,
        functional coverage, a scoreboard and a transaction-level interface,
        with no need to leave \VHDL{} at all. It is mature, well documented,
        and runs under \tool{QuestaSim} and under \tool{NVC}. If your
        organisation is \VHDL{} and intends to stay \VHDL{}, OSVVM is the
        serious answer and deserves a look before a rewrite.
  \item[UVVM, for \VHDL{}.] Also \VHDL{}, also open source, structured around
        \emph{VVCs}, VHDL Verification Components, which occupy the same
        conceptual place as UVM agents, with a central sequencer issuing
        commands to them. UVVM emphasises a directed, readable command
        language over constrained-random generation, which suits control-path
        blocks well.
\end{description}

\begin{ruleofthumb}
Choose the methodology that matches the size of the problem and the size of
the team, not the one on the recruitment advertisement. Then learn UVM
anyway, because the recruitment advertisement is real.
\end{ruleofthumb}

\section{Checklist}
\label{sec:uvm:checklist}

\begin{itemize}
  \item UVM is a \SV{} class library, not a language extension. Every rust
        identifier in this chapter's listings is something you could have
        written yourself, and in \cref{ch:tb} you did.
  \item The six boxes are the same as in \cref{ch:tb}: transaction, stimulus,
        driver, monitor, scoreboard, coverage. UVM adds a base class for
        each, plus the agent, the test class, the factory, the configuration
        database and phasing.
  \item \code{uvm\_component} = structure, built once, has a parent and a
        hierarchical name. \code{uvm\_object} = data, created and destroyed
        constantly, has a name but no place.
  \item Always call \code{super.build\_phase(phase)} first. Omitting it
        breaks configuration and the agent's active/passive switch, silently.
  \item Build is top-down (the parent creates the children); connect is
        bottom-up (both ends must exist). \code{run\_phase} is the only phase
        in which time passes.
  \item Create components and sequence items with
        \code{type\_id::create()}, never with \kw{new}. \kw{new} bypasses the
        factory and every override becomes a no-op with no warning.
  \item Set factory overrides \emph{before} the component is created, which
        means before \code{super.build\_phase()} in the test.
  \item Get the virtual interface from the module into the classes with
        a \code{uvm\_config\_db \#(virtual cordic\_if)::set} call placed
        before \code{run\_test()}, and always test the return value of
        the matching \code{get} with a \chFiveM{uvm\_fatal}.
  \item Raise an objection in the test before starting the sequence and drop
        it after. No objection, no test. Set a drain time long enough for the
        \DUT{} pipeline to empty, and a global timeout so that a stuck
        objection fails instead of hanging.
  \item The driver loop is \code{get\_next\_item(req)}, drive,
        \code{item\_done()}. Missing \code{item\_done()} hangs after exactly
        one transaction.
  \item Randomise the item between \code{start\_item} and \code{finish\_item},
        not before.
  \item Analysis ports broadcast and never block. Coverage collectors get a
        \code{uvm\_subscriber}; scoreboards get a
        \code{uvm\_tlm\_analysis\_fifo}, because checking may take time and
        \code{write()} may not.
  \item Report with \chFiveM{uvm\_info}/\chFiveM{uvm\_warning}/%
        \chFiveM{uvm\_error}/\chFiveM{uvm\_fatal}, never with \code{\$display}
        or \code{\$error}. The verdict is the error count in the report
        summary, and \code{\$error} is not counted.
  \item Make the scoreboard raise an error when it has checked nothing.
        ``Zero failures'' and ``zero checks'' print the same verdict otherwise.
  \item Never call \code{\$finish} yourself: it skips the report summary.
  \item Select the test with \code{+UVM\_TESTNAME}, the log level with
        \code{+UVM\_VERBOSITY}, and always run with a random seed
        (\code{-sv\_seed random}); record the seed so a failure can be
        reproduced.
  \item Debug switches worth memorising: \code{+UVM\_OBJECTION\_TRACE} for a
        test that will not end, \code{+UVM\_CONFIG\_DB\_TRACE} for a value
        that does not arrive, \code{uvm\_top.print\_topology()} for a tree
        that is not what you think it is.
  \item Field automation macros are convenient and slow. Learn with them,
        then write \code{do\_copy}, \code{do\_compare} and
        \code{convert2string} by hand when it matters. Never put a \kw{real}
        in a \chFiveM{uvm\_field\_int}.
  \item Do not reach for UVM on a small block, a short project, or a
        one-person effort. Reach for it when something will be reused: across
        blocks, across levels of hierarchy, across projects, or across a
        team.
\end{itemize}
